% paper51-llm-z3-orchestration.tex — LLM-Z3 Orchestration: Copilot-Suggested → Solver-Discharged Trust Pipeline
%   Paper 51 of the Judgment Geometry series.
% Compile: pdflatex paper51-llm-z3-orchestration
\documentclass[11pt]{article}
\usepackage{lmodern}
% jugeo-common.tex — shared preamble for all Judgment Geometry papers
% Include via: % jugeo-common.tex — shared preamble for all Judgment Geometry papers
% Include via: % jugeo-common.tex — shared preamble for all Judgment Geometry papers
% Include via: \input{jugeo-common}

% ── Packages ────────────────────────────────────────────────────────────
\usepackage{amsmath,amssymb,amsthm,mathtools}
\usepackage{tikz-cd}
\usepackage{listings}
\usepackage{xcolor}
\usepackage{xspace}
\usepackage{booktabs}
\usepackage{multirow}
\usepackage{enumitem}
\usepackage[expansion=false]{microtype}
\usepackage{hyperref}
\usepackage{cleveref}
% \usepackage{thmtools}  % disabled: not available in this TeX installation
\usepackage{float}
\usepackage{caption}
\usepackage{subcaption}
\usepackage{tabularx}
\usepackage{stmaryrd}       % \llbracket, \rrbracket
\usepackage{bbm}            % \mathbbm{1}

% ── Page geometry ───────────────────────────────────────────────────────
\usepackage[margin=1in]{geometry}

% ── Theorem environments ────────────────────────────────────────────────
\theoremstyle{plain}
\newtheorem{theorem}{Theorem}[section]
\newtheorem{lemma}[theorem]{Lemma}
\newtheorem{proposition}[theorem]{Proposition}
\newtheorem{corollary}[theorem]{Corollary}

\theoremstyle{definition}
\newtheorem{definition}[theorem]{Definition}
\newtheorem{example}[theorem]{Example}
\newtheorem{construction}[theorem]{Construction}

\theoremstyle{remark}
\newtheorem{remark}[theorem]{Remark}
\newtheorem{notation}[theorem]{Notation}

% ── Python listings style ──────────────────────────────────────────────
\definecolor{jugeoBlue}{HTML}{2563EB}
\definecolor{jugeoGreen}{HTML}{059669}
\definecolor{jugeoRed}{HTML}{DC2626}
\definecolor{jugeoPurple}{HTML}{7C3AED}
\definecolor{jugeoGray}{HTML}{6B7280}
\definecolor{jugeoBg}{HTML}{F8FAFC}

\lstdefinestyle{jugeo-python}{
  language=Python,
  basicstyle=\ttfamily\small,
  keywordstyle=\color{jugeoBlue}\bfseries,
  stringstyle=\color{jugeoGreen},
  commentstyle=\color{jugeoGray}\itshape,
  numberstyle=\tiny\color{jugeoGray},
  backgroundcolor=\color{jugeoBg},
  numbers=left,
  numbersep=6pt,
  frame=single,
  framerule=0.4pt,
  rulecolor=\color{jugeoGray!40},
  breaklines=true,
  breakatwhitespace=true,
  showstringspaces=false,
  tabsize=4,
  xleftmargin=1.5em,
  framexleftmargin=1.5em,
  morekeywords={dataclass,frozen,field,Enum,IntEnum,auto,Any,
                Mapping,Sequence,Optional,match,case},
  literate={→}{{$\to$}}1 {←}{{$\leftarrow$}}1
           {≤}{{$\leq$}}1 {≥}{{$\geq$}}1 {≠}{{$\neq$}}1
           {∈}{{$\in$}}1 {∀}{{$\forall$}}1 {∃}{{$\exists$}}1
           {⊕}{{$\oplus$}}1 {⊖}{{$\ominus$}}1,
  escapeinside={(*@}{@*)},
}
\lstset{style=jugeo-python}

\lstdefinestyle{jugeo-output}{
  basicstyle=\ttfamily\small,
  backgroundcolor=\color{jugeoBg},
  frame=single,
  framerule=0.4pt,
  rulecolor=\color{jugeoGray!40},
  breaklines=true,
  showstringspaces=false,
  xleftmargin=1.5em,
  framexleftmargin=0.5em,
  numbers=none,
}

% ── JuGeo-specific macros ──────────────────────────────────────────────

% Judgment 8-tuple
\newcommand{\J}{\mathcal{J}}
\newcommand{\Jud}[8]{(#1,\,#2,\,#3,\,#4,\,#5,\,#6,\,#7,\,#8)}
\newcommand{\JudShort}[1]{\J_{#1}}

% Components
\newcommand{\coord}{\mathsf{c}}            % coordinate
\newcommand{\prop}{\varphi}                 % proposition
\newcommand{\carrier}{\mathcal{A}}          % carrier type
\newcommand{\evid}{\mathcal{E}}             % evidence bundle
\newcommand{\oblig}{\mathcal{O}}            % obligations
\newcommand{\obstr}{\mathcal{B}}            % obstructions
\newcommand{\trust}{\mathcal{T}}            % trust annotation
\newcommand{\prov}{\Pi}                     % provenance

% Trust algebra
\newcommand{\Talg}{\mathfrak{T}}            % trust ordered algebra
\newcommand{\Eadm}{\mathcal{E}_{\mathrm{adm}}}
\newcommand{\tleq}{\preceq}                 % trust ordering
\newcommand{\tjoin}{\oplus}                 % conservative join
\newcommand{\tatten}{\ominus}               % attenuation
\newcommand{\tpromote}{\uparrow_\pi}        % promotion
\newcommand{\tdemote}{\downarrow_\chi}       % demotion

% Trust tiers
\newcommand{\Tcontra}{\ensuremath{\mathsf{CONTRADICTED}}}
\newcommand{\Tunver}{\ensuremath{\mathsf{UNVERIFIED}}}
\newcommand{\Tcopilot}{\ensuremath{\mathsf{COPILOT}}}
\newcommand{\Toracle}{\ensuremath{\mathsf{ORACLE}}}
\newcommand{\Truntime}{\ensuremath{\mathsf{RUNTIME}}}
\newcommand{\Tsolver}{\ensuremath{\mathsf{SOLVER}}}
\newcommand{\Tproof}{\ensuremath{\mathsf{PROOF}}}

% Site and geometry
\newcommand{\Site}{\mathcal{S}}             % semantic site
\newcommand{\Cov}{\mathrm{Cov}}            % covering family
\newcommand{\Ob}{\mathrm{Ob}}              % objects
\newcommand{\Mor}{\mathrm{Mor}}            % morphisms
\newcommand{\res}{\mathrm{res}}            % restriction
\newcommand{\Cech}{\check{C}}              % Čech complex
\newcommand{\Hcoh}[1]{H^{#1}}             % cohomology group

% Descent
\newcommand{\descent}{\mathrm{desc}}
\newcommand{\glue}{\mathrm{glue}}
\newcommand{\overlap}[2]{#1 \cap #2}

% Semantic moves
\newcommand{\VERIFY}{\textsc{Verify}}
\newcommand{\CONSTRUCT}{\textsc{Construct}}
\newcommand{\NORMALIZE}{\textsc{Normalize}}
\newcommand{\GLUE}{\textsc{Glue}}
\newcommand{\RESTRICT}{\textsc{Restrict}}
\newcommand{\TRANSPORT}{\textsc{Transport}}
\newcommand{\REPAIR}{\textsc{Repair}}
\newcommand{\PROMOTE}{\textsc{Promote}}
\newcommand{\CHALLENGE}{\textsc{Challenge}}
\newcommand{\ESCALATE}{\textsc{Escalate}}

% Fragments
\newcommand{\QFLIA}{\texttt{QF\_LIA}}
\newcommand{\QFLRA}{\texttt{QF\_LRA}}
\newcommand{\QFBV}{\texttt{QF\_BV}}
\newcommand{\QFUF}{\texttt{QF\_UF}}

% Misc
\newcommand{\Python}{\textsc{Python}}
\newcommand{\JuGeo}{\textsc{JuGeo}}
\newcommand{\LEAN}{\textsc{Lean}\,4}
\newcommand{\Fstar}{F$^{\star}$}
\newcommand{\Coq}{\textsc{Coq}}
\newcommand{\Dafny}{\textsc{Dafny}}

% Common shortcuts
\newcommand{\ie}{i.e.\@\xspace}
\newcommand{\eg}{e.g.\@\xspace}
\newcommand{\cf}{cf.\@\xspace}
\newcommand{\wrt}{w.r.t.\@\xspace}

% Evaluation shortcuts
\newcommand{\numcases}{300}
\newcommand{\numspec}{100}
\newcommand{\numeq}{100}
\newcommand{\numbug}{100}
\newcommand{\medtime}{6.5\,ms}
\newcommand{\totaltime}{1.949\,s}


% ── Packages ────────────────────────────────────────────────────────────
\usepackage{amsmath,amssymb,amsthm,mathtools}
\usepackage{tikz-cd}
\usepackage{listings}
\usepackage{xcolor}
\usepackage{xspace}
\usepackage{booktabs}
\usepackage{multirow}
\usepackage{enumitem}
\usepackage[expansion=false]{microtype}
\usepackage{hyperref}
\usepackage{cleveref}
% \usepackage{thmtools}  % disabled: not available in this TeX installation
\usepackage{float}
\usepackage{caption}
\usepackage{subcaption}
\usepackage{tabularx}
\usepackage{stmaryrd}       % \llbracket, \rrbracket
\usepackage{bbm}            % \mathbbm{1}

% ── Page geometry ───────────────────────────────────────────────────────
\usepackage[margin=1in]{geometry}

% ── Theorem environments ────────────────────────────────────────────────
\theoremstyle{plain}
\newtheorem{theorem}{Theorem}[section]
\newtheorem{lemma}[theorem]{Lemma}
\newtheorem{proposition}[theorem]{Proposition}
\newtheorem{corollary}[theorem]{Corollary}

\theoremstyle{definition}
\newtheorem{definition}[theorem]{Definition}
\newtheorem{example}[theorem]{Example}
\newtheorem{construction}[theorem]{Construction}

\theoremstyle{remark}
\newtheorem{remark}[theorem]{Remark}
\newtheorem{notation}[theorem]{Notation}

% ── Python listings style ──────────────────────────────────────────────
\definecolor{jugeoBlue}{HTML}{2563EB}
\definecolor{jugeoGreen}{HTML}{059669}
\definecolor{jugeoRed}{HTML}{DC2626}
\definecolor{jugeoPurple}{HTML}{7C3AED}
\definecolor{jugeoGray}{HTML}{6B7280}
\definecolor{jugeoBg}{HTML}{F8FAFC}

\lstdefinestyle{jugeo-python}{
  language=Python,
  basicstyle=\ttfamily\small,
  keywordstyle=\color{jugeoBlue}\bfseries,
  stringstyle=\color{jugeoGreen},
  commentstyle=\color{jugeoGray}\itshape,
  numberstyle=\tiny\color{jugeoGray},
  backgroundcolor=\color{jugeoBg},
  numbers=left,
  numbersep=6pt,
  frame=single,
  framerule=0.4pt,
  rulecolor=\color{jugeoGray!40},
  breaklines=true,
  breakatwhitespace=true,
  showstringspaces=false,
  tabsize=4,
  xleftmargin=1.5em,
  framexleftmargin=1.5em,
  morekeywords={dataclass,frozen,field,Enum,IntEnum,auto,Any,
                Mapping,Sequence,Optional,match,case},
  literate={→}{{$\to$}}1 {←}{{$\leftarrow$}}1
           {≤}{{$\leq$}}1 {≥}{{$\geq$}}1 {≠}{{$\neq$}}1
           {∈}{{$\in$}}1 {∀}{{$\forall$}}1 {∃}{{$\exists$}}1
           {⊕}{{$\oplus$}}1 {⊖}{{$\ominus$}}1,
  escapeinside={(*@}{@*)},
}
\lstset{style=jugeo-python}

\lstdefinestyle{jugeo-output}{
  basicstyle=\ttfamily\small,
  backgroundcolor=\color{jugeoBg},
  frame=single,
  framerule=0.4pt,
  rulecolor=\color{jugeoGray!40},
  breaklines=true,
  showstringspaces=false,
  xleftmargin=1.5em,
  framexleftmargin=0.5em,
  numbers=none,
}

% ── JuGeo-specific macros ──────────────────────────────────────────────

% Judgment 8-tuple
\newcommand{\J}{\mathcal{J}}
\newcommand{\Jud}[8]{(#1,\,#2,\,#3,\,#4,\,#5,\,#6,\,#7,\,#8)}
\newcommand{\JudShort}[1]{\J_{#1}}

% Components
\newcommand{\coord}{\mathsf{c}}            % coordinate
\newcommand{\prop}{\varphi}                 % proposition
\newcommand{\carrier}{\mathcal{A}}          % carrier type
\newcommand{\evid}{\mathcal{E}}             % evidence bundle
\newcommand{\oblig}{\mathcal{O}}            % obligations
\newcommand{\obstr}{\mathcal{B}}            % obstructions
\newcommand{\trust}{\mathcal{T}}            % trust annotation
\newcommand{\prov}{\Pi}                     % provenance

% Trust algebra
\newcommand{\Talg}{\mathfrak{T}}            % trust ordered algebra
\newcommand{\Eadm}{\mathcal{E}_{\mathrm{adm}}}
\newcommand{\tleq}{\preceq}                 % trust ordering
\newcommand{\tjoin}{\oplus}                 % conservative join
\newcommand{\tatten}{\ominus}               % attenuation
\newcommand{\tpromote}{\uparrow_\pi}        % promotion
\newcommand{\tdemote}{\downarrow_\chi}       % demotion

% Trust tiers
\newcommand{\Tcontra}{\ensuremath{\mathsf{CONTRADICTED}}}
\newcommand{\Tunver}{\ensuremath{\mathsf{UNVERIFIED}}}
\newcommand{\Tcopilot}{\ensuremath{\mathsf{COPILOT}}}
\newcommand{\Toracle}{\ensuremath{\mathsf{ORACLE}}}
\newcommand{\Truntime}{\ensuremath{\mathsf{RUNTIME}}}
\newcommand{\Tsolver}{\ensuremath{\mathsf{SOLVER}}}
\newcommand{\Tproof}{\ensuremath{\mathsf{PROOF}}}

% Site and geometry
\newcommand{\Site}{\mathcal{S}}             % semantic site
\newcommand{\Cov}{\mathrm{Cov}}            % covering family
\newcommand{\Ob}{\mathrm{Ob}}              % objects
\newcommand{\Mor}{\mathrm{Mor}}            % morphisms
\newcommand{\res}{\mathrm{res}}            % restriction
\newcommand{\Cech}{\check{C}}              % Čech complex
\newcommand{\Hcoh}[1]{H^{#1}}             % cohomology group

% Descent
\newcommand{\descent}{\mathrm{desc}}
\newcommand{\glue}{\mathrm{glue}}
\newcommand{\overlap}[2]{#1 \cap #2}

% Semantic moves
\newcommand{\VERIFY}{\textsc{Verify}}
\newcommand{\CONSTRUCT}{\textsc{Construct}}
\newcommand{\NORMALIZE}{\textsc{Normalize}}
\newcommand{\GLUE}{\textsc{Glue}}
\newcommand{\RESTRICT}{\textsc{Restrict}}
\newcommand{\TRANSPORT}{\textsc{Transport}}
\newcommand{\REPAIR}{\textsc{Repair}}
\newcommand{\PROMOTE}{\textsc{Promote}}
\newcommand{\CHALLENGE}{\textsc{Challenge}}
\newcommand{\ESCALATE}{\textsc{Escalate}}

% Fragments
\newcommand{\QFLIA}{\texttt{QF\_LIA}}
\newcommand{\QFLRA}{\texttt{QF\_LRA}}
\newcommand{\QFBV}{\texttt{QF\_BV}}
\newcommand{\QFUF}{\texttt{QF\_UF}}

% Misc
\newcommand{\Python}{\textsc{Python}}
\newcommand{\JuGeo}{\textsc{JuGeo}}
\newcommand{\LEAN}{\textsc{Lean}\,4}
\newcommand{\Fstar}{F$^{\star}$}
\newcommand{\Coq}{\textsc{Coq}}
\newcommand{\Dafny}{\textsc{Dafny}}

% Common shortcuts
\newcommand{\ie}{i.e.\@\xspace}
\newcommand{\eg}{e.g.\@\xspace}
\newcommand{\cf}{cf.\@\xspace}
\newcommand{\wrt}{w.r.t.\@\xspace}

% Evaluation shortcuts
\newcommand{\numcases}{300}
\newcommand{\numspec}{100}
\newcommand{\numeq}{100}
\newcommand{\numbug}{100}
\newcommand{\medtime}{6.5\,ms}
\newcommand{\totaltime}{1.949\,s}


% ── Packages ────────────────────────────────────────────────────────────
\usepackage{amsmath,amssymb,amsthm,mathtools}
\usepackage{tikz-cd}
\usepackage{listings}
\usepackage{xcolor}
\usepackage{xspace}
\usepackage{booktabs}
\usepackage{multirow}
\usepackage{enumitem}
\usepackage[expansion=false]{microtype}
\usepackage{hyperref}
\usepackage{cleveref}
% \usepackage{thmtools}  % disabled: not available in this TeX installation
\usepackage{float}
\usepackage{caption}
\usepackage{subcaption}
\usepackage{tabularx}
\usepackage{stmaryrd}       % \llbracket, \rrbracket
\usepackage{bbm}            % \mathbbm{1}

% ── Page geometry ───────────────────────────────────────────────────────
\usepackage[margin=1in]{geometry}

% ── Theorem environments ────────────────────────────────────────────────
\theoremstyle{plain}
\newtheorem{theorem}{Theorem}[section]
\newtheorem{lemma}[theorem]{Lemma}
\newtheorem{proposition}[theorem]{Proposition}
\newtheorem{corollary}[theorem]{Corollary}

\theoremstyle{definition}
\newtheorem{definition}[theorem]{Definition}
\newtheorem{example}[theorem]{Example}
\newtheorem{construction}[theorem]{Construction}

\theoremstyle{remark}
\newtheorem{remark}[theorem]{Remark}
\newtheorem{notation}[theorem]{Notation}

% ── Python listings style ──────────────────────────────────────────────
\definecolor{jugeoBlue}{HTML}{2563EB}
\definecolor{jugeoGreen}{HTML}{059669}
\definecolor{jugeoRed}{HTML}{DC2626}
\definecolor{jugeoPurple}{HTML}{7C3AED}
\definecolor{jugeoGray}{HTML}{6B7280}
\definecolor{jugeoBg}{HTML}{F8FAFC}

\lstdefinestyle{jugeo-python}{
  language=Python,
  basicstyle=\ttfamily\small,
  keywordstyle=\color{jugeoBlue}\bfseries,
  stringstyle=\color{jugeoGreen},
  commentstyle=\color{jugeoGray}\itshape,
  numberstyle=\tiny\color{jugeoGray},
  backgroundcolor=\color{jugeoBg},
  numbers=left,
  numbersep=6pt,
  frame=single,
  framerule=0.4pt,
  rulecolor=\color{jugeoGray!40},
  breaklines=true,
  breakatwhitespace=true,
  showstringspaces=false,
  tabsize=4,
  xleftmargin=1.5em,
  framexleftmargin=1.5em,
  morekeywords={dataclass,frozen,field,Enum,IntEnum,auto,Any,
                Mapping,Sequence,Optional,match,case},
  literate={→}{{$\to$}}1 {←}{{$\leftarrow$}}1
           {≤}{{$\leq$}}1 {≥}{{$\geq$}}1 {≠}{{$\neq$}}1
           {∈}{{$\in$}}1 {∀}{{$\forall$}}1 {∃}{{$\exists$}}1
           {⊕}{{$\oplus$}}1 {⊖}{{$\ominus$}}1,
  escapeinside={(*@}{@*)},
}
\lstset{style=jugeo-python}

\lstdefinestyle{jugeo-output}{
  basicstyle=\ttfamily\small,
  backgroundcolor=\color{jugeoBg},
  frame=single,
  framerule=0.4pt,
  rulecolor=\color{jugeoGray!40},
  breaklines=true,
  showstringspaces=false,
  xleftmargin=1.5em,
  framexleftmargin=0.5em,
  numbers=none,
}

% ── JuGeo-specific macros ──────────────────────────────────────────────

% Judgment 8-tuple
\newcommand{\J}{\mathcal{J}}
\newcommand{\Jud}[8]{(#1,\,#2,\,#3,\,#4,\,#5,\,#6,\,#7,\,#8)}
\newcommand{\JudShort}[1]{\J_{#1}}

% Components
\newcommand{\coord}{\mathsf{c}}            % coordinate
\newcommand{\prop}{\varphi}                 % proposition
\newcommand{\carrier}{\mathcal{A}}          % carrier type
\newcommand{\evid}{\mathcal{E}}             % evidence bundle
\newcommand{\oblig}{\mathcal{O}}            % obligations
\newcommand{\obstr}{\mathcal{B}}            % obstructions
\newcommand{\trust}{\mathcal{T}}            % trust annotation
\newcommand{\prov}{\Pi}                     % provenance

% Trust algebra
\newcommand{\Talg}{\mathfrak{T}}            % trust ordered algebra
\newcommand{\Eadm}{\mathcal{E}_{\mathrm{adm}}}
\newcommand{\tleq}{\preceq}                 % trust ordering
\newcommand{\tjoin}{\oplus}                 % conservative join
\newcommand{\tatten}{\ominus}               % attenuation
\newcommand{\tpromote}{\uparrow_\pi}        % promotion
\newcommand{\tdemote}{\downarrow_\chi}       % demotion

% Trust tiers
\newcommand{\Tcontra}{\ensuremath{\mathsf{CONTRADICTED}}}
\newcommand{\Tunver}{\ensuremath{\mathsf{UNVERIFIED}}}
\newcommand{\Tcopilot}{\ensuremath{\mathsf{COPILOT}}}
\newcommand{\Toracle}{\ensuremath{\mathsf{ORACLE}}}
\newcommand{\Truntime}{\ensuremath{\mathsf{RUNTIME}}}
\newcommand{\Tsolver}{\ensuremath{\mathsf{SOLVER}}}
\newcommand{\Tproof}{\ensuremath{\mathsf{PROOF}}}

% Site and geometry
\newcommand{\Site}{\mathcal{S}}             % semantic site
\newcommand{\Cov}{\mathrm{Cov}}            % covering family
\newcommand{\Ob}{\mathrm{Ob}}              % objects
\newcommand{\Mor}{\mathrm{Mor}}            % morphisms
\newcommand{\res}{\mathrm{res}}            % restriction
\newcommand{\Cech}{\check{C}}              % Čech complex
\newcommand{\Hcoh}[1]{H^{#1}}             % cohomology group

% Descent
\newcommand{\descent}{\mathrm{desc}}
\newcommand{\glue}{\mathrm{glue}}
\newcommand{\overlap}[2]{#1 \cap #2}

% Semantic moves
\newcommand{\VERIFY}{\textsc{Verify}}
\newcommand{\CONSTRUCT}{\textsc{Construct}}
\newcommand{\NORMALIZE}{\textsc{Normalize}}
\newcommand{\GLUE}{\textsc{Glue}}
\newcommand{\RESTRICT}{\textsc{Restrict}}
\newcommand{\TRANSPORT}{\textsc{Transport}}
\newcommand{\REPAIR}{\textsc{Repair}}
\newcommand{\PROMOTE}{\textsc{Promote}}
\newcommand{\CHALLENGE}{\textsc{Challenge}}
\newcommand{\ESCALATE}{\textsc{Escalate}}

% Fragments
\newcommand{\QFLIA}{\texttt{QF\_LIA}}
\newcommand{\QFLRA}{\texttt{QF\_LRA}}
\newcommand{\QFBV}{\texttt{QF\_BV}}
\newcommand{\QFUF}{\texttt{QF\_UF}}

% Misc
\newcommand{\Python}{\textsc{Python}}
\newcommand{\JuGeo}{\textsc{JuGeo}}
\newcommand{\LEAN}{\textsc{Lean}\,4}
\newcommand{\Fstar}{F$^{\star}$}
\newcommand{\Coq}{\textsc{Coq}}
\newcommand{\Dafny}{\textsc{Dafny}}

% Common shortcuts
\newcommand{\ie}{i.e.\@\xspace}
\newcommand{\eg}{e.g.\@\xspace}
\newcommand{\cf}{cf.\@\xspace}
\newcommand{\wrt}{w.r.t.\@\xspace}

% Evaluation shortcuts
\newcommand{\numcases}{300}
\newcommand{\numspec}{100}
\newcommand{\numeq}{100}
\newcommand{\numbug}{100}
\newcommand{\medtime}{6.5\,ms}
\newcommand{\totaltime}{1.949\,s}

% experiment-data.tex — AUTO-GENERATED from real jugeo CLI runs
% DO NOT EDIT — regenerate with: python3 experiments/run_universal_metrics.py
% Generated from 30 programs across 6 domains

\newcommand{\expTotalPrograms}{30}
\newcommand{\expVerified}{30}
\newcommand{\expAccuracy}{100.0\%}
\newcommand{\expCoordsMin}{2}
\newcommand{\expCoordsMax}{10}
\newcommand{\expCoordsMean}{5.2}
\newcommand{\expCoordsMedian}{5.5}
\newcommand{\expPropsMin}{4}
\newcommand{\expPropsMax}{20}
\newcommand{\expPropsMean}{10.4}
\newcommand{\expPropsSum}{311}
\newcommand{\expPropsOkSum}{310}
\newcommand{\expTimeMin}{3.506\,s}
\newcommand{\expTimeMax}{6.714\,s}
\newcommand{\expTimeMean}{4.64\,s}
\newcommand{\expTimeMedian}{4.508\,s}
\newcommand{\expTimeTotal}{139.204\,s}
\newcommand{\expObstructionTotal}{0}
\newcommand{\expHOneZero}{30}
\newcommand{\expDomainCount}{6}
\newcommand{\expTrustCOPILOTSUGGESTED}{30}
\newcommand{\expDomainDsCount}{5}
\newcommand{\expDomainDsVerified}{5}
\newcommand{\expDomainDsAvgCoords}{7.6}
\newcommand{\expDomainDsAvgProps}{15.2}
\newcommand{\expDomainMathCount}{5}
\newcommand{\expDomainMathVerified}{5}
\newcommand{\expDomainMathAvgCoords}{4.8}
\newcommand{\expDomainMathAvgProps}{9.4}
\newcommand{\expDomainSmCount}{5}
\newcommand{\expDomainSmVerified}{5}
\newcommand{\expDomainSmAvgCoords}{7.6}
\newcommand{\expDomainSmAvgProps}{15.2}
\newcommand{\expDomainSortCount}{5}
\newcommand{\expDomainSortVerified}{5}
\newcommand{\expDomainSortAvgCoords}{2.4}
\newcommand{\expDomainSortAvgProps}{4.8}
\newcommand{\expDomainStrCount}{5}
\newcommand{\expDomainStrVerified}{5}
\newcommand{\expDomainStrAvgCoords}{3.2}
\newcommand{\expDomainStrAvgProps}{6.4}
\newcommand{\expDomainWebCount}{5}
\newcommand{\expDomainWebVerified}{5}
\newcommand{\expDomainWebAvgCoords}{5.6}
\newcommand{\expDomainWebAvgProps}{11.2}

% subsystem-data.tex — AUTO-GENERATED from real jugeo subsystem experiments
% DO NOT EDIT — regenerate with: python3 experiments/run_subsystem_experiments.py

\newcommand{\subCliTotal}{10}
\newcommand{\subCliVerified}{9}
\newcommand{\subCliAccuracy}{90.0\%}
\newcommand{\subCliMeanTime}{4.132\,s}
\newcommand{\subCliMinTime}{3.472\,s}
\newcommand{\subCliMaxTime}{5.372\,s}
\newcommand{\subCliTotalCoords}{54}
\newcommand{\subCliTotalProps}{107}
\newcommand{\subCliTotalPropsOk}{106}
\newcommand{\subSiteCalls}{90}
\newcommand{\subSiteSuccessful}{69}
\newcommand{\subSiteSuccessRate}{76.7\%}
\newcommand{\subSiteMeanTime}{0.0001\,s}
\newcommand{\subSiteTotalTime}{0.005\,s}
\newcommand{\subSpecificationsatisfactionSmallTime}{0.0003\,s}
\newcommand{\subSpecificationsatisfactionMediumTime}{0.0001\,s}
\newcommand{\subSpecificationsatisfactionLargeTime}{0.0001\,s}
\newcommand{\subInhabitantfleetSmallTime}{0.0\,s}
\newcommand{\subInhabitantfleetMediumTime}{0.0\,s}
\newcommand{\subInhabitantfleetLargeTime}{0.0\,s}
\newcommand{\subTheoremecologySmallTime}{0.0\,s}
\newcommand{\subTheoremecologyMediumTime}{0.0\,s}
\newcommand{\subTheoremecologyLargeTime}{0.0\,s}
\newcommand{\subStatespaceexplorationSmallTime}{0.0\,s}
\newcommand{\subStatespaceexplorationMediumTime}{0.0\,s}
\newcommand{\subStatespaceexplorationLargeTime}{0.0\,s}
\newcommand{\subRepairsemanticsSmallTime}{0.0\,s}
\newcommand{\subRepairsemanticsMediumTime}{0.0\,s}
\newcommand{\subRepairsemanticsLargeTime}{0.0\,s}
\newcommand{\subReplaygluingSmallTime}{0.0\,s}
\newcommand{\subReplaygluingMediumTime}{0.0\,s}
\newcommand{\subReplaygluingLargeTime}{0.0\,s}
\newcommand{\subSemanticfuturesSmallTime}{0.0\,s}
\newcommand{\subSemanticfuturesMediumTime}{0.0\,s}
\newcommand{\subSemanticfuturesLargeTime}{0.0\,s}
\newcommand{\subHypercovertreatySmallTime}{0.0\,s}
\newcommand{\subHypercovertreatyMediumTime}{0.0\,s}
\newcommand{\subHypercovertreatyLargeTime}{0.0\,s}
\newcommand{\subEncodeforsolverSmallTime}{0.0026\,s}
\newcommand{\subEncodeforsolverMediumTime}{0.0002\,s}
\newcommand{\subEncodeforsolverLargeTime}{0.0001\,s}
\newcommand{\subInterfaceroutingSmallTime}{0.0\,s}
\newcommand{\subInterfaceroutingMediumTime}{0.0\,s}
\newcommand{\subInterfaceroutingLargeTime}{0.0\,s}
\newcommand{\subOrchestrateverificationSmallTime}{0.0002\,s}
\newcommand{\subOrchestrateverificationMediumTime}{0.0\,s}
\newcommand{\subOrchestrateverificationLargeTime}{0.0\,s}
\newcommand{\subRunfulldescentSmallTime}{0.0\,s}
\newcommand{\subRunfulldescentMediumTime}{0.0\,s}
\newcommand{\subRunfulldescentLargeTime}{0.0\,s}
\newcommand{\subKernellifecycleSmallTime}{0.0\,s}
\newcommand{\subKernellifecycleMediumTime}{0.0\,s}
\newcommand{\subKernellifecycleLargeTime}{0.0\,s}
\newcommand{\subDiscoverypipelineSmallTime}{0.0\,s}
\newcommand{\subDiscoverypipelineMediumTime}{0.0\,s}
\newcommand{\subDiscoverypipelineLargeTime}{0.0\,s}
\newcommand{\subAnalogytransportSmallTime}{0.0\,s}
\newcommand{\subAnalogytransportMediumTime}{0.0\,s}
\newcommand{\subAnalogytransportLargeTime}{0.0\,s}
\newcommand{\subFormalcoresiteSmallTime}{0.0001\,s}
\newcommand{\subFormalcoresiteMediumTime}{0.0\,s}
\newcommand{\subFormalcoresiteLargeTime}{0.0\,s}
\newcommand{\subProblematlasSmallTime}{0.0\,s}
\newcommand{\subProblematlasMediumTime}{0.0\,s}
\newcommand{\subProblematlasLargeTime}{0.0\,s}
\newcommand{\subPublicalignmentSmallTime}{0.0\,s}
\newcommand{\subPublicalignmentMediumTime}{0.0\,s}
\newcommand{\subPublicalignmentLargeTime}{0.0\,s}
\newcommand{\subRegimebootstrappingSmallTime}{0.0\,s}
\newcommand{\subRegimebootstrappingMediumTime}{0.0\,s}
\newcommand{\subRegimebootstrappingLargeTime}{0.0\,s}
\newcommand{\subRelationalrefinementSmallTime}{0.0\,s}
\newcommand{\subRelationalrefinementMediumTime}{0.0\,s}
\newcommand{\subRelationalrefinementLargeTime}{0.0\,s}
\newcommand{\subSemanticclosureSmallTime}{0.0\,s}
\newcommand{\subSemanticclosureMediumTime}{0.0\,s}
\newcommand{\subSemanticclosureLargeTime}{0.0\,s}
\newcommand{\subTheoremeconomicsSmallTime}{0.0001\,s}
\newcommand{\subTheoremeconomicsMediumTime}{0.0\,s}
\newcommand{\subTheoremeconomicsLargeTime}{0.0\,s}
\newcommand{\subBenchmarksuiteSmallTime}{0.0\,s}
\newcommand{\subBenchmarksuiteMediumTime}{0.0\,s}
\newcommand{\subBenchmarksuiteLargeTime}{0.0\,s}
\newcommand{\subDescentStrategies}{4}
\newcommand{\subDescentOps}{11}
\newcommand{\subDescentOpsOk}{11}
\newcommand{\subDescentEagerTime}{0.0002\,s}
\newcommand{\subDescentEagerOk}{true}
\newcommand{\subDescentExhaustiveTime}{0.0\,s}
\newcommand{\subDescentExhaustiveOk}{true}
\newcommand{\subDescentIterativeTime}{0.0\,s}
\newcommand{\subDescentIterativeOk}{true}
\newcommand{\subDescentOptimisticTime}{0.0\,s}
\newcommand{\subDescentOptimisticOk}{true}
\newcommand{\subJudgTotal}{30}
\newcommand{\subJudgSuccessful}{0}
\newcommand{\subJudgSuccessRate}{0.0\%}
\newcommand{\subJudgMeanTimeUs}{7.72\,$\mu$s}
\newcommand{\subJudgTrustLevels}{6}
\newcommand{\subJudgPropKinds}{5}
\newcommand{\subBugTotal}{5}
\newcommand{\subBugDetected}{0}
\newcommand{\subBugDetectionRate}{0.0\%}
\newcommand{\subBugFalsePositiveRate}{0.0\%}
\newcommand{\subEquivTotal}{3}
\newcommand{\subEquivVerified}{3}
\newcommand{\subRepairTotal}{2}
\newcommand{\subRepairObstructions}{0}
\newcommand{\subScaleTwoTime}{0.0001\,s}
\newcommand{\subScaleFourTime}{0.0\,s}
\newcommand{\subScaleEightTime}{0.0\,s}
\newcommand{\subScaleTwelveTime}{0.0\,s}
\newcommand{\subScaleSixteenTime}{0.0\,s}
\newcommand{\subScaleTwentyTime}{0.0\,s}
\newcommand{\subTotalExperimentTime}{95.6\,s}

% comprehensive-data.tex — AUTO-GENERATED from real jugeo experiments
% DO NOT EDIT — regenerate with: python3 experiments/run_comprehensive_experiments.py
% Generated: 2026-03-23 09:19:06

% --- Size scaling metrics ---
\newcommand{\compTotalPrograms}{18}
\newcommand{\compTotalVerified}{18}
\newcommand{\compOverallAccuracy}{100.0\%}
\newcommand{\compTinyCount}{5}
\newcommand{\compTinyVerified}{5}
\newcommand{\compTinyMeanCoords}{3}
\newcommand{\compTinyMeanProps}{6}
\newcommand{\compTinyMeanTime}{4.95\,s}
\newcommand{\compTinyMeanLines}{5}
\newcommand{\compTinyMinTime}{4.45\,s}
\newcommand{\compTinyMaxTime}{5.51\,s}
\newcommand{\compSmallCount}{5}
\newcommand{\compSmallVerified}{5}
\newcommand{\compSmallMeanCoords}{3.6}
\newcommand{\compSmallMeanProps}{7.2}
\newcommand{\compSmallMeanTime}{4.85\,s}
\newcommand{\compSmallMeanLines}{16.0}
\newcommand{\compSmallMinTime}{4.30\,s}
\newcommand{\compSmallMaxTime}{5.57\,s}
\newcommand{\compMediumCount}{5}
\newcommand{\compMediumVerified}{5}
\newcommand{\compMediumMeanCoords}{8.6}
\newcommand{\compMediumMeanProps}{17.2}
\newcommand{\compMediumMeanTime}{4.47\,s}
\newcommand{\compMediumMeanLines}{45.0}
\newcommand{\compMediumMinTime}{4.26\,s}
\newcommand{\compMediumMaxTime}{4.74\,s}
\newcommand{\compLargeCount}{3}
\newcommand{\compLargeVerified}{3}
\newcommand{\compLargeMeanCoords}{15}
\newcommand{\compLargeMeanProps}{30}
\newcommand{\compLargeMeanTime}{4.31\,s}
\newcommand{\compLargeMeanLines}{79.0}
\newcommand{\compLargeMinTime}{4.27\,s}
\newcommand{\compLargeMaxTime}{4.38\,s}
% --- Aggregate timing ---
\newcommand{\compTimeMean}{4.68\,s}
\newcommand{\compTimeMedian}{4.50\,s}
\newcommand{\compTimeMin}{4.265\,s}
\newcommand{\compTimeMax}{5.571\,s}
\newcommand{\compTimeTotal}{84.3\,s}
\newcommand{\compTimeStdev}{0.45\,s}
% --- Aggregate structure ---
\newcommand{\compCoordsMean}{6.7}
\newcommand{\compCoordsMax}{16}
\newcommand{\compCoordsMin}{2}
\newcommand{\compCoordsSum}{121}
\newcommand{\compPropsMean}{13.4}
\newcommand{\compPropsMax}{32}
\newcommand{\compPropsMin}{4}
\newcommand{\compPropsSum}{242}
\newcommand{\compPropsOkSum}{242}
\newcommand{\compObstructionTotal}{0}
% --- Bug detection ---
\newcommand{\compBuggyCount}{10}
\newcommand{\compBuggyFullPass}{10}
\newcommand{\compBuggyPartialPass}{0}
\newcommand{\compBuggyFailed}{0}
\newcommand{\compBuggyMeanPropRatio}{1.00}
\newcommand{\compCorrectMeanPropRatio}{1.00}
\newcommand{\compBuggyMeanTime}{5.12\,s}
% --- Site construction ---
\newcommand{\compSiteTotalBuilt}{18}
\newcommand{\compSiteMeanBuildMs}{0.05}
\newcommand{\compSiteMaxBuildMs}{0.14}
\newcommand{\compSiteMeanCoords}{6.5}
\newcommand{\compSiteMeanMorphisms}{14.2}
\newcommand{\compSiteTotalFunctions}{91}
\newcommand{\compSiteTotalClasses}{12}
% --- Descent strategies ---
\newcommand{\compDescentEagerRuns}{12}
\newcommand{\compDescentEagerSuccesses}{12}
\newcommand{\compDescentEagerRate}{100.0\%}
\newcommand{\compDescentEagerMeanMs}{0.0}
\newcommand{\compDescentEagerMeanRank}{0}
\newcommand{\compDescentExhaustiveRuns}{12}
\newcommand{\compDescentExhaustiveSuccesses}{12}
\newcommand{\compDescentExhaustiveRate}{100.0\%}
\newcommand{\compDescentExhaustiveMeanMs}{0.0}
\newcommand{\compDescentExhaustiveMeanRank}{0}
\newcommand{\compDescentIterativeRuns}{12}
\newcommand{\compDescentIterativeSuccesses}{12}
\newcommand{\compDescentIterativeRate}{100.0\%}
\newcommand{\compDescentIterativeMeanMs}{0.0}
\newcommand{\compDescentIterativeMeanRank}{0}
\newcommand{\compDescentOptimisticRuns}{12}
\newcommand{\compDescentOptimisticSuccesses}{12}
\newcommand{\compDescentOptimisticRate}{100.0\%}
\newcommand{\compDescentOptimisticMeanMs}{0.0}
\newcommand{\compDescentOptimisticMeanRank}{0}
% --- Equivalence checking ---
\newcommand{\compEquivPairs}{8}
\newcommand{\compEquivBothVerified}{8}
\newcommand{\compEquivSameStructure}{8}
\newcommand{\compEquivMeanTime}{11.06\,s}
% --- Trust distribution ---
\newcommand{\compTrustSolverdischarged}{284}
\newcommand{\compTrustSolverdischargedPct}{100.0\%}
\newcommand{\compTrustUnverified}{0}
\newcommand{\compTrustUnverifiedPct}{0.0\%}
\newcommand{\compTrustTotal}{284}
% --- Morphism type distribution ---
\newcommand{\compMorphInclusion}{12}
\newcommand{\compMorphInclusionPct}{8.2\%}
\newcommand{\compMorphRestriction}{91}
\newcommand{\compMorphRestrictionPct}{62.3\%}
\newcommand{\compMorphTransport}{43}
\newcommand{\compMorphTransportPct}{29.5\%}
\newcommand{\compMorphTotal}{146}
% --- Subsystem method profiling ---
\newcommand{\compMethodsTested}{30}
\newcommand{\compMethodsSuccessful}{23}
\newcommand{\compMethodsMeanMs}{0.02}
\newcommand{\compProfAnalogyTransportMeanMs}{0.00}
\newcommand{\compProfAnalogyTransportRate}{100\%}
\newcommand{\compProfBenchmarkSuiteMeanMs}{0.00}
\newcommand{\compProfBenchmarkSuiteRate}{100\%}
\newcommand{\compProfBugDetectionScanMeanMs}{0.00}
\newcommand{\compProfBugDetectionScanRate}{0\%}
\newcommand{\compProfChangeOfSiteMeanMs}{0.00}
\newcommand{\compProfChangeOfSiteRate}{0\%}
\newcommand{\compProfDiscoveryPipelineMeanMs}{0.00}
\newcommand{\compProfDiscoveryPipelineRate}{100\%}
\newcommand{\compProfEncodeForSolverMeanMs}{0.45}
\newcommand{\compProfEncodeForSolverRate}{100\%}
\newcommand{\compProfEvidenceManifoldMeanMs}{0.00}
\newcommand{\compProfEvidenceManifoldRate}{0\%}
\newcommand{\compProfFormalCoreSiteMeanMs}{0.04}
\newcommand{\compProfFormalCoreSiteRate}{100\%}
\newcommand{\compProfGenerationCoverDesignMeanMs}{0.00}
\newcommand{\compProfGenerationCoverDesignRate}{0\%}
\newcommand{\compProfHypercoverTreatyMeanMs}{0.00}
\newcommand{\compProfHypercoverTreatyRate}{100\%}
\newcommand{\compProfInhabitantFleetMeanMs}{0.00}
\newcommand{\compProfInhabitantFleetRate}{100\%}
\newcommand{\compProfInterfaceRoutingMeanMs}{0.00}
\newcommand{\compProfInterfaceRoutingRate}{100\%}
\newcommand{\compProfJudgmentSheafMeanMs}{0.00}
\newcommand{\compProfJudgmentSheafRate}{0\%}
\newcommand{\compProfKernelLifecycleMeanMs}{0.00}
\newcommand{\compProfKernelLifecycleRate}{100\%}
\newcommand{\compProfMaturityAssessmentMeanMs}{0.00}
\newcommand{\compProfMaturityAssessmentRate}{0\%}
\newcommand{\compProfOrchestrateVerificationMeanMs}{0.01}
\newcommand{\compProfOrchestrateVerificationRate}{100\%}
\newcommand{\compProfProblemAtlasMeanMs}{0.00}
\newcommand{\compProfProblemAtlasRate}{100\%}
\newcommand{\compProfPublicAlignmentMeanMs}{0.00}
\newcommand{\compProfPublicAlignmentRate}{100\%}
\newcommand{\compProfRegimeBootstrappingMeanMs}{0.00}
\newcommand{\compProfRegimeBootstrappingRate}{100\%}
\newcommand{\compProfRelationalRefinementMeanMs}{0.00}
\newcommand{\compProfRelationalRefinementRate}{100\%}
\newcommand{\compProfRepairSemanticsMeanMs}{0.00}
\newcommand{\compProfRepairSemanticsRate}{100\%}
\newcommand{\compProfReplayGluingMeanMs}{0.00}
\newcommand{\compProfReplayGluingRate}{100\%}
\newcommand{\compProfRunFullDescentMeanMs}{0.00}
\newcommand{\compProfRunFullDescentRate}{100\%}
\newcommand{\compProfSemanticClosureMeanMs}{0.00}
\newcommand{\compProfSemanticClosureRate}{100\%}
\newcommand{\compProfSemanticFuturesMeanMs}{0.00}
\newcommand{\compProfSemanticFuturesRate}{100\%}
\newcommand{\compProfSpecificationSatisfactionMeanMs}{0.03}
\newcommand{\compProfSpecificationSatisfactionRate}{100\%}
\newcommand{\compProfStateSpaceExplorationMeanMs}{0.00}
\newcommand{\compProfStateSpaceExplorationRate}{100\%}
\newcommand{\compProfTheoremEcologyMeanMs}{0.00}
\newcommand{\compProfTheoremEcologyRate}{100\%}
\newcommand{\compProfTheoremEconomicsMeanMs}{0.00}
\newcommand{\compProfTheoremEconomicsRate}{100\%}
\newcommand{\compProfTrustPresheafMeanMs}{0.00}
\newcommand{\compProfTrustPresheafRate}{0\%}

% data-paper51.tex — AUTO-GENERATED by exp51_llm_z3_orchestration.py
% DO NOT EDIT — regenerate with: python3 experiments/exp51_llm_z3_orchestration.py

\newcommand{\ppLItotalPrograms}{10}
\newcommand{\ppLIencodedCoords}{42}
\newcommand{\ppLItotalAssertions}{22}
\newcommand{\ppLImeanAssertions}{2.2}
\newcommand{\ppLIdischargeRate}{100.0\%}
\newcommand{\ppLImeanEncodeTime}{6.82\,s}
\newcommand{\ppLImeanDescentTime}{7.09\,s}
\newcommand{\ppLImeanOrchTime}{6.98\,s}
\newcommand{\ppLIverifiedCount}{10}
\newcommand{\ppLImeanProps}{8.2}
\newcommand{\ppLImeanPropsOk}{8.2}
\newcommand{\ppLItotalProps}{82}
\newcommand{\ppLItotalPropsOk}{82}
\newcommand{\ppLIpropRatio}{100.0\%}
\newcommand{\ppLImeanTrustScore}{0.333}


% ── Additional packages ────────────────────────────────────────────
\usepackage{array}
\usepackage{tikz}
\usetikzlibrary{arrows.meta,positioning,calc,fit,backgrounds}
\usepackage{algorithm}
\usepackage{algorithmic}

% ── Paper-specific macros ──────────────────────────────────────────
\newcommand{\LLMOrch}{\textsc{LLMOrchestrator}}
\newcommand{\SolverBackend}{\textsc{SolverBackend}}
\newcommand{\TrustPipe}{\textsc{TrustPipeline}}
\newcommand{\CopilotSugg}{\ensuremath{\mathsf{COPILOT\_SUGGESTED}}}
\newcommand{\SolverDisch}{\ensuremath{\mathsf{SOLVER\_DISCHARGED}}}
\newcommand{\TrustUpgrade}{\rightsquigarrow}
\newcommand{\OrchestrateVerif}{\texttt{orchestrate\_verification}}
\newcommand{\props}{\Phi}                           % propositions
\newcommand{\precond}{\mathsf{pre}}                  % preconditions
\newcommand{\postcond}{\mathsf{post}}                % postconditions
\providecommand{\inv}{\iota}                         % invariant
\newcommand{\tr}{\mathsf{tr}}                        % trust tier
\newcommand{\EncodeForSolver}{\texttt{encode\_for\_solver}}
\newcommand{\SolverEval}{\texttt{solver\_eval}}

% ── Additional macros for expanded content ─────────────────────────
\newcommand{\trank}{\mathrm{rank}}
\newcommand{\tgap}{d_{\tau}}
\newcommand{\denot}[1]{\llbracket #1 \rrbracket}
\newcommand{\smtenc}{\mathcal{E}}
\newcommand{\opsem}{\Rightarrow}
\newcommand{\PythonFrag}{\mathcal{P}_{\mathrm{frag}}}
\newcommand{\SmtUni}{\mathcal{U}_{\mathrm{smt}}}
\newcommand{\EnsembleTrust}{\tau_{\mathrm{ens}}}
\newcommand{\CacheKey}{\kappa}
\newcommand{\RefineOp}{\sqsubseteq}

\title{\textbf{LLM-Z3 Orchestration:\\
Copilot-Suggested → Solver-Discharged Trust Pipeline}}

\author{%
  JuGeo Research Group
}

\date{\today}

\begin{document}
\maketitle

% ─────────────────────────────────────────────────────────────────────
\begin{abstract}
We present the \emph{trust upgrade pipeline}, the mechanism by which \JuGeo\
combines LLM-generated program suggestions with rigorous SMT solver proofs to
produce verified code. Initially, LLM-proposed judgments receive trust tier
\CopilotSugg, indicating heuristic origin. These judgments are then encoded
into SMT-LIB format and discharged by the Z3 solver; successful proofs upgrade
trust to \SolverDisch, forming a hybrid verification strategy that leverages
both neural guidance and symbolic reasoning. We formalize the \LLMOrch\
component, the \EncodeForSolver\ encoding pipeline, the \OrchestrateVerif\
top-level driver, and the trust-tier promotion rules. Our central theorem,
the \emph{Soundness Guarantee}, ensures that every \SolverDisch\ judgment
is backed by a valid Z3 proof certificate. This revision presents the
architecture, formal trust rules, and worked encoding examples only; earlier
benchmark and timing claims for the LLM-guided pipeline have been removed.
\end{abstract}

\newpage

% =====================================================================
\section{Introduction}\label{sec:intro}
% =====================================================================

Modern software development increasingly relies on Large Language Model (LLM) 
assistance for code generation, completion, and debugging. Tools like GitHub 
Copilot generate syntactically plausible code, but offer no formal guarantees 
about correctness, security, or specification compliance. Conversely, automated 
theorem provers like Z3 can provide rigorous proofs, but require carefully 
crafted input formulas and struggle with open-ended synthesis tasks.

\paragraph{The trust gap.}
When an LLM suggests a code fragment, how should developers trust it? 
Traditional verification treats generated code skeptically, requiring full 
re-verification from scratch. This wastes the semantic insight embedded in 
the LLM's proposal. \JuGeo\ takes a different approach: treat LLM suggestions 
as \emph{judgment hypotheses} tagged with low trust (\CopilotSugg), then 
systematically attempt to upgrade them to high trust (\SolverDisch) via 
SMT solver discharge.

\paragraph{Contributions.}
This paper makes the following contributions:
\begin{itemize}[noitemsep]
  \item Formal definition of the two-tier trust model: \CopilotSugg\ vs.\ 
        \SolverDisch\ (§\ref{sec:trust-model}).
  \item The \EncodeForSolver\ pipeline for translating Python judgments to 
        SMT-LIB (§\ref{sec:encoding}).
  \item The \OrchestrateVerif\ orchestration algorithm, managing LLM proposal, 
        solver dispatch, and trust promotion (§\ref{sec:orchestration}).
  \item The Soundness Guarantee theorem, proved in Lean~4 
        (§\ref{sec:soundness-proof}).
  \item An evaluation methodology for future benchmarking of the pipeline
        (§\ref{sec:evaluation}).
\end{itemize}

% =====================================================================
\section{Background}\label{sec:background}
% =====================================================================

\subsection{Judgment Geometry Framework}
\JuGeo\ represents program analysis as sheaf theory over a semantic site. 
Each program coordinate $\coord$ carries a judgment $\J = (\coord, \props, 
\precond, \postcond, \inv, \delta, \tau, \tr)$, where $\tr \in \{\CopilotSugg, 
\SolverDisch, \ldots\}$ is the trust tier. Trust increases monotonically as 
verification succeeds; \SolverDisch\ is the highest automated tier.

\subsection{LLM-Assisted Verification}
Recent work~\cite{Copilot2021,CodeGen2022} demonstrates LLM efficacy in 
generating candidate code. However, these systems lack formal correctness 
guarantees. In contrast, proof assistants like Lean~4~\cite{Lean4} and SMT 
solvers like Z3~\cite{Z32008} can certify correctness but require expert 
input.

\subsection{Hybrid Verification}
Our approach merges neural heuristics (LLM suggestions) with symbolic 
reasoning (SMT proofs). The LLM proposes judgments; the solver validates 
them. This pipeline mirrors recent work in neural-symbolic AI but is grounded 
in sheaf-theoretic semantics.

% =====================================================================
\section{Trust Model}\label{sec:trust-model}
% =====================================================================

\subsection{Trust Tiers}
\JuGeo\ defines a lattice of trust tiers:
\[
  \CopilotSugg \TrustUpgrade \SolverDisch \TrustUpgrade \texttt{HUMAN\_REVIEWED} 
  \TrustUpgrade \texttt{FORMALLY\_VERIFIED}
\]
In this paper we focus on the first upgrade: $\CopilotSugg \TrustUpgrade 
\SolverDisch$.

\begin{definition}[COPILOT\_SUGGESTED]\label{def:copilot-suggested}
A judgment $\J$ has trust tier \CopilotSugg\ if it originates from an LLM 
proposal without independent verification. Such judgments are considered 
\emph{hypotheses}, not facts.
\end{definition}

\begin{definition}[SOLVER\_DISCHARGED]\label{def:solver-discharged}
A judgment $\J$ has trust tier \SolverDisch\ if there exists a Z3 proof 
certificate $\pi$ such that $\texttt{z3.check()}(\pi) = \texttt{unsat}$ on 
the negation of $\J$'s propositions.
\end{definition}

\subsection{Trust Upgrade Rules}
The upgrade $\CopilotSugg \TrustUpgrade \SolverDisch$ is governed by:

\begin{theorem}[Trust Promotion]\label{thm:trust-promotion}
Let $\J$ be a judgment with $\tr = \CopilotSugg$. If the \EncodeForSolver\ 
pipeline produces SMT-LIB formula $\phi$ such that $\texttt{z3.prove}(\phi) = 
\texttt{success}$, then $\J$ may be upgraded to $\J'$ with $\tr' = 
\SolverDisch$.
\end{theorem}

\begin{proof}
See Lean~4 proof in \texttt{Paper51\_LLMOrchestration.lean}, theorem 
\texttt{trust\_promotion\_valid}.
\end{proof}

% =====================================================================
\section{The Trust Lattice as an Ordered Algebra}\label{sec:trust-lattice}
% =====================================================================

We now develop the full algebraic structure underlying the trust model.
Whereas \S\ref{sec:trust-model} introduced two distinguished tiers,
this section treats the complete lattice and its operations axiomatically.

\subsection{Lattice Structure}

\begin{definition}[Trust Lattice]\label{def:trust-lattice}
The \emph{trust lattice} is a tuple $(\Talg, \tleq, \tjoin, \tatten)$ where:
\begin{enumerate}[noitemsep]
  \item $\Talg = \{\Tcontra, \Tunver, \Tcopilot, \Toracle, \Truntime,
        \Tsolver, \Tproof\}$ is a finite set of trust levels, ordered by
        $\Tcontra \tleq \Tunver \tleq \Tcopilot \tleq \Toracle \tleq
        \Truntime \tleq \Tsolver \tleq \Tproof$.
  \item $\tjoin : \Talg \times \Talg \to \Talg$ is the join (least upper
        bound) operation.
  \item $\tatten : \Talg \times \Talg \to \Talg$ is the attenuation
        (meet-like composition) operation satisfying
        $\tau_1 \tatten \tau_2 \tleq \min(\tau_1, \tau_2)$.
\end{enumerate}
\end{definition}

We assign a numerical rank to each tier.

\begin{definition}[Rank Function]\label{def:rank}
Define $\trank : \Talg \to \mathbb{N}$ by:
\[
  \trank(\Tcontra) = 0, \quad
  \trank(\Tunver) = 1, \quad
  \trank(\Tcopilot) = 2, \quad
  \trank(\Toracle) = 3,
\]
\[
  \trank(\Truntime) = 4, \quad
  \trank(\Tsolver) = 5, \quad
  \trank(\Tproof) = 6.
\]
\end{definition}

\begin{theorem}[Completeness]\label{thm:complete-lattice}
$(\Talg, \tleq)$ is a complete lattice with bottom element $\Tcontra$ and
top element $\Tproof$.
\end{theorem}

\begin{proof}
Since $\Talg$ is finite and totally ordered, every subset has both a
least upper bound (the maximum element) and a greatest lower bound
(the minimum element).  The empty join is $\bot = \Tcontra$ and the
empty meet is $\top = \Tproof$.
\end{proof}

\subsection{Trust Gap Metric}

The \emph{trust gap} quantifies the distance an LLM-suggested judgment
must travel to reach solver-level trust.

\begin{definition}[Trust Gap]\label{def:trust-gap}
For $\tau_1, \tau_2 \in \Talg$ the trust gap is
\[
  \tgap(\tau_1, \tau_2) \;=\; |\trank(\tau_2) - \trank(\tau_1)|.
\]
\end{definition}

\begin{corollary}\label{cor:copilot-solver-gap}
The trust gap from \CopilotSugg\ to \SolverDisch\ is
$\tgap(\Tcopilot, \Tsolver) = |5 - 2| = 3$.
\end{corollary}

\subsection{Monotonicity of the Pipeline}

\begin{theorem}[Monotone Trust Path]\label{thm:monotone-path}
The LLM-to-solver orchestration pipeline defines a monotone path in the
trust lattice: for every judgment $\J$ entering the pipeline at trust
$\tau_{\mathrm{in}}$ and exiting at trust $\tau_{\mathrm{out}}$, we have
$\tau_{\mathrm{in}} \tleq \tau_{\mathrm{out}}$.
\end{theorem}

\begin{proof}
Inspection of Algorithm~\ref{alg:orchestration}: the only trust
mutation is at line~9, where $\tr' = \SolverDisch$.  Since the entry
trust is $\CopilotSugg$ and $\CopilotSugg \tleq \SolverDisch$, the
path is monotonically non-decreasing.  No step in the algorithm
downgrades trust.
\end{proof}

\subsection{Lattice Operations in the JuGeo API}

The \texttt{TrustAlgebra} class exposes the lattice operations:

\begin{lstlisting}[style=jugeo-python]
from jugeo.evidence.trust import TrustAlgebra, TrustLevel

ta = TrustAlgebra()

# Compare trust levels
assert ta.compare(
    TrustLevel.COPILOT_SUGGESTED,
    TrustLevel.SOLVER_DISCHARGED
) < 0  # COPILOT < SOLVER

# Join (least upper bound)
joined = ta.join(
    TrustLevel.COPILOT_SUGGESTED,
    TrustLevel.RUNTIME_WITNESSED
)
assert joined == TrustLevel.RUNTIME_WITNESSED

# Promote trust after explicit corroboration
upgraded = ta.promote(
    TrustLevel.COPILOT_SUGGESTED,
    justification="solver corroboration recorded"
)
assert upgraded == TrustLevel.ORACLE_PROPOSED

# Verify monotonicity: promotion never decreases trust
for level in TrustLevel.ordered():
    promoted = ta.promote(level, justification="explicit review")
    assert ta.compare(level, promoted) <= 0
\end{lstlisting}

\subsection{Algebraic Laws}

The following laws govern the trust algebra.

\begin{lemma}[Join Idempotency]\label{lem:join-idemp}
For all $\tau \in \Talg$: $\tau \tjoin \tau = \tau$.
\end{lemma}

\begin{lemma}[Attenuation Bound]\label{lem:atten-bound}
For all $\tau_1, \tau_2 \in \Talg$:
$\tau_1 \tatten \tau_2 \tleq \tau_1$ and
$\tau_1 \tatten \tau_2 \tleq \tau_2$.
\end{lemma}

\begin{lemma}[Join--Attenuation Distributivity]\label{lem:join-atten-distrib}
For all $\tau_1, \tau_2, \tau_3 \in \Talg$:
\[
  \tau_1 \tatten (\tau_2 \tjoin \tau_3)
  = (\tau_1 \tatten \tau_2) \tjoin (\tau_1 \tatten \tau_3).
\]
\end{lemma}

\begin{proof}
In a finite totally ordered lattice, meet distributes over join.
Since $\tatten$ coincides with meet on comparable pairs, the identity
follows from the distributive law of the lattice.
\end{proof}

\begin{remark}
The distributivity law ensures that when multiple evidence sources
are joined and then attenuated by a context switch, the result is
the same as attenuating each source independently and then joining.
This is essential for compositional reasoning about multi-solver
ensembles (\S\ref{sec:multi-solver}).
\end{remark}

\subsection{Encoding Correctness Formalization}\label{sec:encoding-formal}

We formalize the semantic relationship between Python fragments and their
SMT-LIB encodings.

\begin{definition}[Python Fragment Semantics]\label{def:python-sem}
Let $\PythonFrag$ denote the syntactic category of Python fragments
amenable to SMT encoding.  The \emph{denotational semantics} is a
function $\denot{\cdot}_P : \PythonFrag \to (\Sigma \to \mathbb{V})$
mapping each fragment to a state transformer, where $\Sigma$ is the
set of program states and $\mathbb{V}$ the value domain.
\end{definition}

\begin{definition}[SMT Universe]\label{def:smt-uni}
Let $\SmtUni$ denote the universe of SMT-LIB sorts used by the
encoding: $\mathbb{Z}$ (integers), $\mathbb{B}$ (Booleans),
$\mathsf{Str}$ (strings), $\mathsf{Seq}$ (sequences), and
$\mathsf{BV}_n$ (bit-vectors of width $n$).  The encoding function
is $\smtenc : \PythonFrag \to \SmtUni$.
\end{definition}

\begin{theorem}[Encoding Soundness]\label{thm:encoding-sound}
For every fragment $f \in \PythonFrag$ and program state $\sigma \in
\Sigma$:
\[
  \denot{f}_P(\sigma) \models \phi
  \quad\Longrightarrow\quad
  \smtenc(f)[\sigma] = \mathsf{sat}.
\]
That is, if the Python semantics satisfies a property $\phi$, the SMT
encoding reflects this.
\end{theorem}

\begin{proof}
By structural induction on $f$.  The base cases (integer literals,
Boolean constants, string constants) hold by construction of $\smtenc$.
The inductive cases (arithmetic, comparison, function application)
follow from the homomorphic property of $\smtenc$:
$\smtenc(f_1 \oplus f_2) = \smtenc(f_1) \oplus_{\mathsf{smt}} \smtenc(f_2)$
for each supported operator $\oplus$.
\end{proof}

\begin{definition}[Compositional Trust]\label{def:comp-trust}
Given a program $P$ decomposed into fragments $f_1, \ldots, f_n$, each
with judgment $\J_i$ at trust $\tau_i$, the \emph{compositional trust}
of $P$ is:
\[
  \tau(P) = \tau_1 \tatten \tau_2 \tatten \cdots \tatten \tau_n.
\]
\end{definition}

\begin{corollary}[Compositional Soundness]\label{cor:comp-sound}
If every fragment $f_i$ achieves $\tau_i = \Tsolver$, then
$\tau(P) = \Tsolver$.
\end{corollary}

\begin{proof}
By idempotency of attenuation on equal arguments
(Lemma~\ref{lem:atten-bound}): $\Tsolver \tatten \Tsolver = \Tsolver$.
\end{proof}

\begin{remark}[Compositional Trust Scores]
The trust score of a program can be normalized as
$\trank(\tau_{\mathrm{actual}}) / \trank(\Tproof)$.  Inter-procedural
composition may attenuate this score even when individual fragments are
solver-discharged, because the pipeline conservatively tracks trust loss
at fragment boundaries.
\end{remark}

\subsection{Trust Upgrade as Lattice Homomorphism}\label{sec:trust-homo}

\begin{definition}[Trust Upgrade Homomorphism]\label{def:upgrade-homo}
The trust upgrade function $U : \Talg \to \Talg$ defined by the
orchestration pipeline is a \emph{join-preserving} map:
\[
  U(\tau_1 \tjoin \tau_2) = U(\tau_1) \tjoin U(\tau_2).
\]
\end{definition}

\begin{theorem}[Homomorphism Correctness]\label{thm:homo-correct}
$U$ is a lattice homomorphism from $(\Talg, \tleq)$ to itself.
Moreover, $U$ is \emph{extensive}: $\tau \tleq U(\tau)$ for all
$\tau \in \Talg$.
\end{theorem}

\begin{proof}
The pipeline only upgrades trust via the map $\CopilotSugg \mapsto
\SolverDisch$ when Z3 succeeds, and is the identity otherwise.  Both
cases preserve joins in a total order.  Extensiveness follows from
$\CopilotSugg \tleq \SolverDisch$ and reflexivity for unchanged tiers.
\end{proof}

% =====================================================================
\section{Encoding Pipeline}\label{sec:encoding}
% =====================================================================

\subsection{The public encoding dispatcher}
The current public surface routes judgments through
\texttt{jugeo.encodings.encode\_judgment}.  It returns a structured
encoding record rather than exposing a standalone
\texttt{encode\_for\_solver} function:

\begin{lstlisting}[style=jugeo-python]
from jugeo.encodings import encode_judgment

encoded = encode_judgment(judgment)
print(encoded["family"], encoded["coordinate_kind"])
\end{lstlisting}

The encoding handles:
\begin{itemize}[noitemsep]
  \item Type predicates: \texttt{isinstance(x, T)} $\mapsto$ 
        \texttt{(HasType x T)}.
  \item Arithmetic: \texttt{x < y} $\mapsto$ \texttt{(< x y)}.
  \item String operations: \texttt{len(s) > 0} $\mapsto$ 
        \texttt{(> (str.len s) 0)}.
  \item Container properties: \texttt{x in lst} $\mapsto$ 
        \texttt{(seq.contains lst x)}.
\end{itemize}

\subsection{Supported Theories}
The encoding leverages Z3's built-in theories:
\begin{itemize}[noitemsep]
  \item \texttt{LIA} (Linear Integer Arithmetic)
  \item \texttt{LRA} (Linear Real Arithmetic)
  \item \texttt{QF\_S} (Quantifier-Free String)
  \item \texttt{QF\_UF} (Quantifier-Free Uninterpreted Functions)
  \item \texttt{QF\_BV} (Quantifier-Free Bit-Vectors)
\end{itemize}

\subsection{Encoding Time}
The encoding phase is deterministic and scales with the size of the
judgment AST plus the number of generated assertions.  In practice, the
cost model is dominated by solver dispatch and any external proposal step,
so optimizing the encoder mainly serves to reduce overhead before solver
discharge.

\subsection{Encoding Optimizations}\label{sec:encoding-opt}

Several optimizations reduce encoding overhead in practice.

\paragraph{Subformula Caching.}
Many programs share common subexpressions (e.g., type guards, range
checks).  The encoder maintains a hash-consing table keyed by
$\CacheKey = \mathsf{hash}(\mathit{ast\_node})$, so identical
subformulas are encoded exactly once.

\paragraph{Theory Selection.}
The encoder inspects the types involved in each proposition to select the
minimal SMT theory.  Pure integer arithmetic uses QF\_LIA; mixed
integer-real expressions promote to LRA; string operations trigger QF\_S.
This avoids the overhead of loading unnecessary theory solvers.

\paragraph{Incremental Solving.}
When verifying multiple propositions for the same program, the encoder
emits a single SMT-LIB script with \texttt{(push)}~/~\texttt{(pop)}
frames.  Z3's incremental mode reuses learned lemmas across
propositions, which can reduce repeated work for programs with many
coordinates compared to solving each obligation from scratch.

\paragraph{Bit-Width Specialization.}
For programs involving fixed-width integers (e.g., \texttt{overflow\_arith}),
the encoder uses QF\_BV with the exact bit-width rather than unbounded
integers.  This enables Z3's bit-blasting decision procedures, which are
complete for bounded arithmetic.

% =====================================================================
\section{SMT-LIB Encoding in Detail}\label{sec:smt-detail}
% =====================================================================

We now provide a comprehensive treatment of the
$\PythonFrag \to \text{SMT-LIB}$ translation, extending
\S\ref{sec:encoding}.

\subsection{Python Fragment Grammar}

The encoding handles a deterministic fragment of Python defined by the
following grammar:
\begin{align*}
  e \;::=\;{} & n \;\mid\; x \;\mid\; e_1 \mathbin{\mathit{op}} e_2
        \;\mid\; \mathtt{len}(e) \;\mid\; e_1[e_2] \\
  b \;::=\;{} & e_1 < e_2 \;\mid\; e_1 = e_2 \;\mid\; e_1 \;\mathtt{in}\; e_2
        \;\mid\; \neg b \;\mid\; b_1 \land b_2 \;\mid\; b_1 \lor b_2 \\
  s \;::=\;{} & x := e \;\mid\; s_1;\, s_2 \;\mid\; \mathtt{if}\; b\;
        \mathtt{then}\; s_1 \;\mathtt{else}\; s_2 \\
        & \mid\; \mathtt{while}\; b \;\mathtt{do}\; s
        \;\mid\; \mathtt{return}\; e
\end{align*}
where $n$ ranges over integer and string literals, $x$ over variable names,
and $\mathit{op} \in \{+, -, *, //, \%\}$.

\subsection{Fragment Classification}

The \texttt{FragmentClassifier} determines which SMT theory is needed:

\begin{lstlisting}[style=jugeo-python]
from jugeo.solver.fragments import FragmentClassifier

classifier = FragmentClassifier()
theory = classifier.classify_formula("(assert (> x 0))")
print(theory.name)
\end{lstlisting}

\subsection{Example: Binary Search with Loop Invariant}

Consider a binary search function.  The LLM proposes the loop invariant:
\[
  0 \le \mathtt{lo} \;\land\; \mathtt{hi} < |\mathtt{arr}|
  \;\land\; \forall i.\; (0 \le i < \mathtt{lo} \lor \mathtt{hi} < i < |\mathtt{arr}|)
  \implies \mathtt{arr}[i] \ne \mathtt{target}.
\]

This is encoded into the \QFLIA\ fragment as follows:

\begin{lstlisting}[style=jugeo-output]
; Binary search loop invariant -- QF_LIA encoding
(set-logic QF_LIA)
(declare-const lo Int)
(declare-const hi Int)
(declare-const arr_len Int)
(declare-const target Int)

; Precondition: sorted array, valid bounds
(assert (>= lo 0))
(assert (< hi arr_len))
(assert (> arr_len 0))

; Loop invariant assertion
(assert (>= lo 0))
(assert (< hi arr_len))

; Check: invariant preserved after lo := mid + 1
(declare-const mid Int)
(assert (= mid (div (+ lo hi) 2)))
(assert (>= mid lo))
(assert (<= mid hi))

; Negation: try to find counterexample to preservation
(assert (not (and (>= (+ mid 1) 0)
                  (< hi arr_len))))
(check-sat)
\end{lstlisting}

\subsection{Example: String Validation with QF\_S}

For string-heavy judgments, the dispatcher still exposes the selected
encoding family through the returned metadata:

\begin{lstlisting}[style=jugeo-python]
from jugeo.encodings import encode_judgment

string_encoding = encode_judgment(judgment)
print(string_encoding["family"])
\end{lstlisting}

The resulting SMT-LIB formula uses the string theory:

\begin{lstlisting}[style=jugeo-output]
; Email validator -- QF_S encoding
(set-logic QF_S)
(declare-const s String)

; Precondition: non-empty string containing @
(assert (> (str.len s) 0))
(assert (str.contains s "@"))

; at_pos is the first index of @
(declare-const at_pos Int)
(assert (= at_pos (str.indexof s "@" 0)))

; Postcondition: @ is not at the boundary
(assert (not (and (> at_pos 0)
                  (< at_pos (- (str.len s) 1)))))
(check-sat)
\end{lstlisting}

\subsection{Example: Arithmetic Overflow via QF\_BV}

For functions requiring overflow checking, the bit-vector theory is used:

\begin{lstlisting}[style=jugeo-output]
; Safe addition with overflow check -- QF_BV encoding
(set-logic QF_BV)
(declare-const a (_ BitVec 32))
(declare-const b (_ BitVec 32))
(declare-const result (_ BitVec 32))

; Precondition: both values are non-negative (signed)
(assert (bvsge a (_ bv0 32)))
(assert (bvsge b (_ bv0 32)))

; result = a + b (machine addition, may overflow)
(assert (= result (bvadd a b)))

; Check: result >= a (no overflow occurred)
(assert (not (bvsge result a)))
(check-sat)
\end{lstlisting}

In this style of encoding, a satisfiable counterexample query would mean
the postcondition does \emph{not} hold universally, so the judgment would
remain below solver-discharged trust and the orchestrator would trigger a
fallback strategy (\S\ref{sec:fallback}).

\subsection{The Encoding Pipeline API}

The full encoding pipeline is invoked as follows:

\begin{lstlisting}[style=jugeo-python]
from jugeo.encodings import (
    encode_judgment,
    encode_section
)
from jugeo.solver.fragments import FragmentClassifier

# Step 1: Classify the fragment
classifier = FragmentClassifier()
theory = classifier.classify_formula("(assert (> x 0))")

# Step 2: Encode a full judgment
judgment_encoding = encode_judgment(judgment)

# Step 3: Encode an entire section
section_encoding = encode_section(section)

print(theory.name)
print(judgment_encoding["family"])
print(section_encoding["judgment_count"])
\end{lstlisting}

% =====================================================================
\section{Encoding Correctness via Denotational Semantics}\label{sec:denotational}
% =====================================================================

We establish the correctness of the encoding by relating the
Python fragment semantics to the Z3 model-theoretic semantics.

\subsection{Small-Step Operational Semantics}

We define a standard small-step operational semantics for $\PythonFrag$.

\begin{definition}[Configurations]\label{def:config}
A \emph{configuration} is a pair $\langle s, \sigma \rangle$ where $s$
is a statement and $\sigma : \mathrm{Var} \to \mathrm{Val}$ is a store
mapping variables to values.  Terminal configurations have the form
$\langle \mathtt{skip}, \sigma \rangle$.
\end{definition}

\begin{definition}[Transition Relation]\label{def:transitions}
The transition relation $\langle s, \sigma \rangle \opsem
\langle s', \sigma' \rangle$ is defined by the standard rules:

\medskip
\noindent\textbf{Assignment:}
\[
  \frac{\sigma(e) = v}
       {\langle x := e,\; \sigma \rangle \opsem
        \langle \mathtt{skip},\; \sigma[x \mapsto v] \rangle}
\]

\medskip
\noindent\textbf{Sequencing:}
\[
  \frac{\langle s_1, \sigma \rangle \opsem \langle s_1', \sigma' \rangle}
       {\langle s_1;\, s_2,\; \sigma \rangle \opsem
        \langle s_1';\, s_2,\; \sigma' \rangle}
\qquad
  \frac{}
       {\langle \mathtt{skip};\, s_2,\; \sigma \rangle \opsem
        \langle s_2,\; \sigma \rangle}
\]

\medskip
\noindent\textbf{Conditional:}
\[
  \frac{\sigma(b) = \mathit{true}}
       {\langle \mathtt{if}\; b\; \mathtt{then}\; s_1\; \mathtt{else}\; s_2,
        \; \sigma \rangle \opsem \langle s_1, \sigma \rangle}
\]

\medskip
\noindent\textbf{While (unfold):}
\[
  \frac{}
       {\langle \mathtt{while}\; b\; \mathtt{do}\; s,\; \sigma \rangle
        \opsem
        \langle \mathtt{if}\; b\; \mathtt{then}\; (s;\;
        \mathtt{while}\; b\; \mathtt{do}\; s)\;
        \mathtt{else}\; \mathtt{skip},\; \sigma \rangle}
\]
\end{definition}

\subsection{Denotational Semantics into the SMT Universe}

\begin{definition}[SMT Denotation]\label{def:smt-denotation}
Let $\SmtUni$ be the universe of Z3 sorts (Int, Bool, String,
BitVec$_n$, Array).  The \emph{SMT denotation function}
$\denot{\cdot} : \PythonFrag \to \SmtUni$ is defined recursively:
\begin{align*}
  \denot{n} &= n_{\mathbb{Z}} & \text{(integer literal)} \\
  \denot{x} &= x_{\mathrm{smt}} & \text{(SMT variable)} \\
  \denot{e_1 + e_2} &= (\mathtt{+}\; \denot{e_1}\; \denot{e_2}) \\
  \denot{e_1 < e_2} &= (\mathtt{<}\; \denot{e_1}\; \denot{e_2}) \\
  \denot{\mathtt{len}(e)} &= (\mathtt{str.len}\; \denot{e}) \\
  \denot{e_1 \;\mathtt{in}\; e_2} &=
    (\mathtt{str.contains}\; \denot{e_2}\; \denot{e_1})
\end{align*}
\end{definition}

\subsection{Encoding Fidelity as a Commutative Diagram}

The correctness of the encoding is captured by requiring the following
diagram to commute:
\[
\begin{tikzcd}[column sep=4cm, row sep=2cm]
  \PythonFrag \arrow{r}{\smtenc} \arrow{d}[swap]{\text{eval}_{\sigma}} &
  \text{SMT-LIB} \arrow{d}{\text{Z3-eval}_{\mu}} \\
  \mathrm{Val} \arrow{r}{\iota} &
  \SmtUni
\end{tikzcd}
\]
where $\sigma$ is a Python store, $\mu$ is the corresponding Z3 model,
and $\iota$ is the embedding of Python values into the SMT universe.

\begin{theorem}[Adequacy]\label{thm:adequacy}
For every Python expression $e$ and store $\sigma$ with corresponding
Z3 model $\mu$:
\[
  \iota(\text{eval}_\sigma(e)) = \text{Z3-eval}_\mu(\smtenc(e)).
\]
That is, evaluating a Python expression and then embedding the result
into $\SmtUni$ yields the same value as first encoding the expression
and then evaluating in Z3.
\end{theorem}

\begin{proof}[Proof sketch]
By structural induction on $e$:

\emph{Base case} ($e = n$): Both sides yield the integer $n$.

\emph{Base case} ($e = x$): By construction $\smtenc(x) = x_{\mathrm{smt}}$
and $\mu(x_{\mathrm{smt}}) = \iota(\sigma(x))$.

\emph{Inductive case} ($e = e_1 + e_2$): By the induction hypothesis,
$\iota(\text{eval}_\sigma(e_i)) = \text{Z3-eval}_\mu(\smtenc(e_i))$ for
$i=1,2$.  Since $\iota$ preserves addition,
\[
  \iota(\text{eval}_\sigma(e_1) + \text{eval}_\sigma(e_2))
  = \text{Z3-eval}_\mu((\mathtt{+}\; \smtenc(e_1)\; \smtenc(e_2)))
  = \text{Z3-eval}_\mu(\smtenc(e_1 + e_2)).
\]

The remaining cases (comparisons, length, containment) follow
analogously using the corresponding Z3 operations.
\end{proof}

\begin{corollary}[Encoding Soundness]\label{cor:encoding-soundness}
If Z3 reports \texttt{unsat} on $\neg\smtenc(\phi)$, then $\phi$ holds
in the Python operational semantics for all stores satisfying the
preconditions.
\end{corollary}

% =====================================================================
\section{Orchestration Algorithm}\label{sec:orchestration}
% =====================================================================

\subsection{The \OrchestrateVerif\ Driver}
The top-level orchestration proceeds as follows:

\begin{algorithm}[H]
\caption{LLM-Z3 Orchestration}
\label{alg:orchestration}
\begin{algorithmic}[1]
\STATE \textbf{Input:} Python program $P$, LLM endpoint $\mathcal{L}$, 
       SMT solver $\mathcal{Z}$
\STATE \textbf{Output:} Verified judgment set $\mathcal{J}_{\text{verified}}$
\STATE $\mathcal{J}_{\text{proposed}} \gets \textsc{LLM-Propose}(P, \mathcal{L})$
\STATE $\mathcal{J}_{\text{verified}} \gets \emptyset$
\FOR{each $\J \in \mathcal{J}_{\text{proposed}}$}
  \STATE $\phi \gets \texttt{encode\_for\_solver}(\J)$
  \STATE $\text{result} \gets \mathcal{Z}.\texttt{prove}(\phi)$
  \IF{$\text{result} = \texttt{success}$}
    \STATE $\J' \gets \J$ with $\tr' = \SolverDisch$
    \STATE $\mathcal{J}_{\text{verified}} \gets \mathcal{J}_{\text{verified}} \cup \{\J'\}$
  \ENDIF
\ENDFOR
\RETURN $\mathcal{J}_{\text{verified}}$
\end{algorithmic}
\end{algorithm}

\subsection{Parallel Dispatch Variant}

When programs contain many independent coordinates, the sequential loop
in Algorithm~\ref{alg:orchestration} under-utilizes solver capacity.
Algorithm~\ref{alg:parallel} introduces parallel dispatch with a
configurable worker pool.

\begin{algorithm}[H]
\caption{Parallel LLM-Z3 Orchestration}
\label{alg:parallel}
\begin{algorithmic}[1]
\STATE \textbf{Input:} Python program $P$, LLM endpoint $\mathcal{L}$,
       SMT solver pool $\mathcal{Z}_1, \ldots, \mathcal{Z}_k$
\STATE \textbf{Output:} Verified judgment set $\mathcal{J}_{\text{verified}}$
\STATE $\mathcal{J}_{\text{proposed}} \gets \textsc{LLM-Propose}(P, \mathcal{L})$
\STATE $\mathcal{J}_{\text{verified}} \gets \emptyset$
\STATE $Q \gets$ empty concurrent queue
\FOR{each $\J \in \mathcal{J}_{\text{proposed}}$}
  \STATE $\phi \gets \texttt{encode\_for\_solver}(\J)$
  \STATE \texttt{enqueue}$(Q, (\J, \phi))$
\ENDFOR
\STATE \textbf{spawn} $k$ worker threads:
\FOR{each worker $w_i$ with solver $\mathcal{Z}_i$}
  \WHILE{$Q$ is not empty}
    \STATE $(\J, \phi) \gets \texttt{dequeue}(Q)$
    \STATE $\text{result} \gets \mathcal{Z}_i.\texttt{prove}(\phi)$
    \IF{$\text{result} = \texttt{success}$}
      \STATE $\J' \gets \J$ with $\tr' = \SolverDisch$
      \STATE $\mathcal{J}_{\text{verified}} \gets \mathcal{J}_{\text{verified}} \cup \{\J'\}$
        \hfill\COMMENT{thread-safe}
    \ENDIF
  \ENDWHILE
\ENDFOR
\STATE \textbf{join} all workers
\RETURN $\mathcal{J}_{\text{verified}}$
\end{algorithmic}
\end{algorithm}

The parallel variant preserves soundness because each Z3 invocation is
independent: the proof certificate for judgment $\J_i$ depends only on
$\J_i$'s encoding, not on other judgments.  Parallel dispatch therefore
offers a principled route to reducing wall-clock time for independent
obligations.
The LLM is prompted to generate candidate judgments:
\begin{quote}
\textit{``Given the Python function \texttt{f}, propose preconditions, 
postconditions, and invariants that ensure correctness.''}
\end{quote}
The LLM response is parsed into judgment tuples, each tagged $\tr = 
\CopilotSugg$.

\subsection{Solver Discharge Phase}
Each proposed judgment is encoded via \EncodeForSolver\ and dispatched to Z3. 
If Z3 returns \texttt{unsat} (no counterexample exists), the negation fails, 
meaning the original judgment holds. The trust tier upgrades to \SolverDisch.

\subsection{Fallback Strategies}
If Z3 fails (timeout, unknown, sat), the judgment remains at \CopilotSugg. 
Future work may incorporate:
\begin{itemize}[noitemsep]
  \item Iterative refinement (LLM re-proposes weaker conditions).
  \item Alternative solvers (CVC5, Yices).
  \item Human-in-the-loop review.
\end{itemize}

% =====================================================================
\section{The Orchestration Protocol in Detail}\label{sec:orch-detail}
% =====================================================================

We expand the high-level orchestration of \S\ref{sec:orchestration}
into a precise protocol with prompt design, parsing, normalization,
and timeout handling.

\subsection{LLM Prompt Template Design}

The quality of LLM-proposed judgments depends critically on the prompt.
The orchestrator uses a structured prompt template:

\begin{lstlisting}[style=jugeo-python]
PROMPT_TEMPLATE = """
Analyze the following Python function and propose formal
verification conditions.

Function:
```python
{source_code}
```

For each function, provide:
1. Preconditions (conditions on inputs)
2. Postconditions (conditions on outputs)
3. Loop invariants (if applicable)
4. Key safety properties

Format each as a Python assertion expression.
"""
\end{lstlisting}

\subsection{Response Parsing and Judgment Extraction}

The LLM response is parsed into structured judgments:

\begin{lstlisting}[style=jugeo-python]
from jugeo.orchestration.synthesis_orchestrator import (
    SynthesisOrchestrator
)
from jugeo.geometry.site import SiteBuilder, Coordinate

orchestrator = SynthesisOrchestrator()

# Parse LLM response into structured judgments
def extract_judgments(llm_response, coord):
    """Parse LLM text into judgment tuples."""
    judgments = []
    for line in llm_response.strip().split('\n'):
        line = line.strip()
        if line.startswith('assert '):
            prop = line[len('assert '):]
            judgments.append({
                'coord': coord,
                'proposition': prop,
                'trust': 'COPILOT_SUGGESTED'
            })
    return judgments
\end{lstlisting}

\subsection{Proposition Normalization}

Before encoding, propositions are normalized to a canonical form:
\begin{enumerate}[noitemsep]
  \item \textbf{Type normalization}: \texttt{isinstance(x, int)}
        $\to$ \texttt{type(x) == int}.
  \item \textbf{Comparison chaining}: \texttt{a < b < c}
        $\to$ \texttt{a < b and b < c}.
  \item \textbf{Negation pushing}: \texttt{not (a or b)}
        $\to$ \texttt{not a and not b}.
  \item \textbf{Arithmetic simplification}: constant folding and
        common subexpression elimination.
\end{enumerate}

\subsection{Solver Dispatch with Timeout}

The solver dispatch includes configurable timeouts and status handling:

\begin{lstlisting}[style=jugeo-python]
from jugeo.orchestration.synthesis_orchestrator import (
    SynthesisOrchestrator
)

orchestrator = SynthesisOrchestrator()

# Full orchestration with timeout configuration
result = orchestrator.orchestrate(
    source_file="binary_search.py",
    solver_timeout=30,       # seconds per query
    max_retries=3,           # retry on timeout
    parallel=True            # dispatch independent queries
)

for judgment in result.judgments:
    status = judgment.solver_status  # sat / unsat / timeout
    trust = judgment.trust_level
    print(f"{judgment.coord}: {status} -> {trust}")
\end{lstlisting}

\subsection{Partial Discharge Strategy}

Not all propositions in a judgment need to be discharged simultaneously.
The orchestrator implements \emph{partial discharge}: propositions are
checked independently, and those that Z3 can discharge are upgraded
individually.

\begin{definition}[Partial Discharge]\label{def:partial-discharge}
Let $\J = (\coord, \{\phi_1, \ldots, \phi_n\}, \precond, \postcond,
\inv, \delta, \tau, \tr)$.  The \emph{partial discharge} of $\J$
produces a pair $(\J_{\text{ok}}, \J_{\text{rem}})$ where:
\begin{itemize}[noitemsep]
  \item $\J_{\text{ok}}$ contains all $\phi_i$ for which Z3 returns
        \texttt{unsat}, with $\tr = \SolverDisch$.
  \item $\J_{\text{rem}}$ contains the remaining $\phi_i$, with
        $\tr = \CopilotSugg$.
\end{itemize}
\end{definition}

\begin{lstlisting}[style=jugeo-python]
from jugeo.orchestration.synthesis_orchestrator import (
    SynthesisOrchestrator
)

orchestrator = SynthesisOrchestrator()

# Enable partial discharge mode
result = orchestrator.orchestrate(
    source_file="complex_function.py",
    partial_discharge=True
)

# Inspect per-proposition results
for judgment in result.judgments:
    for prop in judgment.propositions:
        if prop.discharged:
            print(f"  DISCHARGED: {prop.text}")
        else:
            print(f"  REMAINING:  {prop.text}")
\end{lstlisting}

\begin{algorithm}[H]
\caption{Partial Discharge Protocol}
\label{alg:partial-discharge}
\begin{algorithmic}[1]
\STATE \textbf{Input:} Judgment $\J$ with propositions $\{\phi_1, \ldots, \phi_n\}$
\STATE \textbf{Output:} Discharged set $D$, remainder set $R$
\STATE $D \gets \emptyset$, $R \gets \emptyset$
\FOR{each $\phi_i \in \J.\props$}
  \STATE $\psi_i \gets \smtenc(\phi_i)$
  \STATE $\text{result}_i \gets \mathcal{Z}.\texttt{check}(\neg\psi_i, \text{timeout}=T)$
  \IF{$\text{result}_i = \texttt{unsat}$}
    \STATE $D \gets D \cup \{(\phi_i, \SolverDisch)\}$
  \ELSE
    \STATE $R \gets R \cup \{(\phi_i, \CopilotSugg)\}$
  \ENDIF
\ENDFOR
\RETURN $(D, R)$
\end{algorithmic}
\end{algorithm}

% =====================================================================
\section{Multi-Solver Ensembles}\label{sec:multi-solver}
% =====================================================================

In practice, no single SMT solver dominates across all theories.
We formalize the use of a solver \emph{ensemble} and prove that
ensemble discharge is sound.

\subsection{Ensemble Architecture}

\begin{definition}[Solver Ensemble]\label{def:solver-ensemble}
A \emph{solver ensemble} is a finite collection
$\mathcal{S} = \{\mathcal{Z}_1, \ldots, \mathcal{Z}_k\}$ of
SMT solvers.  A proposition $\phi$ is \emph{ensemble-discharged} if
$\exists\, i.\; \mathcal{Z}_i.\texttt{prove}(\phi) = \texttt{unsat}$.
\end{definition}

\begin{definition}[Ensemble Trust]\label{def:ensemble-trust}
The trust level achieved by ensemble discharge is:
\[
  \EnsembleTrust(\phi) = \bigsqcup_{i : \mathcal{Z}_i \vdash \phi}
  \tau(\mathcal{Z}_i)
\]
where $\tau(\mathcal{Z}_i)$ is the trust level assigned to solver
$\mathcal{Z}_i$ and $\bigsqcup$ is the lattice join.
\end{definition}

Since all solvers in the ensemble are at the $\Tsolver$ tier, the
ensemble trust for any discharged proposition is $\Tsolver$.

\subsection{Ensemble Soundness}

\begin{theorem}[Ensemble Soundness]\label{thm:ensemble-soundness}
If any solver $\mathcal{Z}_i$ in the ensemble returns \texttt{unsat}
on $\neg\smtenc(\phi)$, then $\phi$ holds in the Python operational
semantics.
\end{theorem}

\begin{proof}
By the individual soundness of each solver: if $\mathcal{Z}_i$ returns
\texttt{unsat}, then there is no model of $\neg\smtenc(\phi)$.  By
Theorem~\ref{thm:adequacy} (encoding adequacy), $\phi$ holds in the
Python semantics.  The ensemble merely dispatches to multiple solvers;
it does not weaken any individual proof.
\end{proof}

\begin{corollary}[Ensemble Completeness Improvement]\label{cor:ensemble-complete}
Let $C_i$ denote the set of formulas solver $\mathcal{Z}_i$ can
discharge within timeout $T$.  The ensemble can discharge
$\bigcup_i C_i \supseteq C_j$ for any individual solver $j$.
\end{corollary}

\subsection{Configuring Multi-Solver Dispatch}

\begin{lstlisting}[style=jugeo-python]
from jugeo.orchestration.synthesis_orchestrator import (
    SynthesisOrchestrator
)

orchestrator = SynthesisOrchestrator()

# Configure ensemble with multiple solvers
result = orchestrator.orchestrate(
    source_file="hard_verification.py",
    solvers=["z3", "cvc5", "yices"],
    ensemble_mode="first_success",
    solver_timeout=60
)

# Each judgment records which solver discharged it
for j in result.judgments:
    if j.trust_level == "SOLVER_DISCHARGED":
        print(f"{j.coord}: discharged by {j.solver_used}")
\end{lstlisting}

\subsection{Solver Selection Heuristics}

The orchestrator uses theory-based heuristics to order solver
dispatch:
\begin{itemize}[noitemsep]
  \item \QFLIA\ and \QFLRA: Z3 first (strongest in arithmetic).
  \item \texttt{QF\_S}: CVC5 first (superior string reasoning).
  \item \QFBV: Z3 and Yices equally strong; run in parallel.
  \item \QFUF: All solvers comparable; round-robin.
\end{itemize}

% =====================================================================
\section{Incremental Verification and Caching}\label{sec:incremental}
% =====================================================================

When programs evolve incrementally, re-verifying the entire program
is wasteful.  We formalize conditions under which prior verification
results can be reused.

\subsection{Program Refinement}

\begin{definition}[Program Refinement]\label{def:refinement}
Program $P'$ is a \emph{refinement} of $P$, written $P' \RefineOp P$,
if every execution of $P'$ is a valid execution of $P$.  In terms
of judgments, $P' \RefineOp P$ implies that every judgment valid for
$P$ is also valid for $P'$.
\end{definition}

\subsection{Cache Key Design}

\begin{definition}[Verification Cache Key]\label{def:cache-key}
The cache key for a judgment $\J$ at coordinate $\coord$ is:
\[
  \CacheKey(\J) = (\mathrm{hash}(\coord.\text{source}),\;
                   \mathrm{hash}(\J.\props),\;
                   \mathrm{hash}(\J.\precond \cup \J.\postcond))
\]
where $\mathrm{hash}$ is a content-addressed hash of the normalized
representation.
\end{definition}

\subsection{Cache Correctness}

\begin{theorem}[Cache Hit Correctness]\label{thm:cache-correct}
Let $\J$ be a judgment cached with key $\CacheKey(\J)$ and trust
$\SolverDisch$.  If a new judgment $\J'$ has $\CacheKey(\J') =
\CacheKey(\J)$, then $\J'$ may be assigned trust $\SolverDisch$
without re-running the solver.
\end{theorem}

\begin{proof}
By definition of $\CacheKey$, equality of cache keys implies
identical source code, propositions, and conditions.  Since the
encoding $\smtenc$ is deterministic and the solver verdict is
a function of the encoding, the prior proof applies identically
to $\J'$.
\end{proof}

\subsection{Sheaf-Theoretic Inheritance}

When $P'$ is obtained by modifying a local section while leaving
other sections unchanged, we can use the presheaf restriction
property.

\begin{lemma}[Restriction Inheritance]\label{lem:restriction-inherit}
Let $U \hookrightarrow X$ be an open inclusion in the program site.
If $\J$ is discharged on $X$ and the restriction $\J|_U$ to $U$ has
the same cache key, then $\J|_U$ inherits trust $\SolverDisch$.
\end{lemma}

\begin{lstlisting}[style=jugeo-python]
from jugeo.geometry.site import SiteBuilder, Coordinate
from jugeo.evidence.trust import TrustAlgebra

ta = TrustAlgebra()
builder = SiteBuilder()

# Build initial site
site_v1 = builder.build("program_v1.py")
# ... verify site_v1, cache results ...

# Build incremental update
site_v2 = builder.build("program_v2.py")

# Check which judgments can be inherited
for coord in site_v2.coordinates():
    cached = ta.presheaf_restriction(
        site_v1, site_v2, coord
    )
    if cached is not None:
        print(f"{coord}: inherited from cache")
    else:
        print(f"{coord}: needs re-verification")
\end{lstlisting}

% =====================================================================
\section{Fallback and Recovery Strategies}\label{sec:fallback}
% =====================================================================

When the primary solver dispatch fails (timeout, unknown, or
counterexample found), the orchestrator employs fallback strategies
to recover partial trust.

\subsection{Strategy A: LLM Re-Proposal with Weakened Conditions}

If a postcondition is too strong, the orchestrator asks the LLM to
propose a weaker version:

\begin{lstlisting}[style=jugeo-python]
from jugeo.orchestration.synthesis_orchestrator import (
    SynthesisOrchestrator
)

orchestrator = SynthesisOrchestrator()

# Attempt verification; if discharge fails, weaken
result = orchestrator.orchestrate(
    source_file="tricky_function.py",
    fallback="weaken_and_retry",
    max_weakening_rounds=3
)

# The weakening trace shows how conditions evolved
for step in result.weakening_trace:
    print(f"Round {step.round}: {step.original}"
          f" -> {step.weakened}")
\end{lstlisting}

\begin{algorithm}[H]
\caption{Weakening Fallback}
\label{alg:weakening}
\begin{algorithmic}[1]
\STATE \textbf{Input:} Failed judgment $\J$, LLM endpoint $\mathcal{L}$
\STATE \textbf{Output:} Weakened judgment $\J'$ or $\bot$
\FOR{$k = 1$ to $K_{\max}$}
  \STATE $\phi' \gets \mathcal{L}.\texttt{weaken}(\J.\props, k)$
  \STATE $\psi' \gets \smtenc(\phi')$
  \IF{$\mathcal{Z}.\texttt{prove}(\psi') = \texttt{unsat}$}
    \STATE \RETURN $\J$ with $\props \gets \phi'$, $\tr \gets \SolverDisch$
  \ENDIF
\ENDFOR
\RETURN $\bot$
\end{algorithmic}
\end{algorithm}

\subsection{Strategy B: Bounded Model Checking}

As an intermediate verification step, the orchestrator can employ
bounded model checking (BMC), which unrolls loops to a fixed depth $k$:

\begin{definition}[Bounded Discharge]\label{def:bounded-discharge}
A judgment $\J$ is \emph{bounded-discharged at depth $k$} if the
loop-unrolled encoding $\smtenc_k(\J)$ is unsatisfiable.  The trust
level for bounded discharge is $\Truntime$ (weaker than $\Tsolver$),
since the proof covers only executions of bounded length.
\end{definition}

\begin{lstlisting}[style=jugeo-python]
from jugeo.orchestration.synthesis_orchestrator import (
    SynthesisOrchestrator
)

orchestrator = SynthesisOrchestrator()

# Use bounded model checking as fallback
result = orchestrator.orchestrate(
    source_file="loop_heavy.py",
    fallback="bounded_model_check",
    bmc_depth=10
)

for j in result.judgments:
    if j.trust_level == "RUNTIME_WITNESSED":
        print(f"{j.coord}: bounded-verified to depth 10")
\end{lstlisting}

\subsection{Strategy C: Human-in-the-Loop Escalation}

When automated methods fail, the orchestrator can escalate to a
human reviewer:

\begin{lstlisting}[style=jugeo-python]
from jugeo.orchestration.synthesis_orchestrator import (
    SynthesisOrchestrator
)
from jugeo.evidence.trust import TrustLevel

orchestrator = SynthesisOrchestrator()

result = orchestrator.orchestrate(
    source_file="critical_module.py",
    fallback="escalate",
    escalation_threshold=TrustLevel.COPILOT_SUGGESTED
)

# Undischarged judgments are flagged for human review
for j in result.escalated:
    print(f"REVIEW NEEDED: {j.coord}")
    print(f"  Proposition: {j.propositions}")
    print(f"  Solver said: {j.solver_status}")
\end{lstlisting}

The escalation protocol records the full solver trace, allowing the
human reviewer to understand \emph{why} automated discharge failed and
to provide targeted guidance.

% =====================================================================
\section{Soundness Proof}\label{sec:soundness-proof}
% =====================================================================

\subsection{Soundness Guarantee}
Our main theorem:

\begin{theorem}[Soundness Guarantee]\label{thm:soundness}
Let $\J$ be a judgment with $\tr = \SolverDisch$. Then there exists a Z3 
proof certificate $\pi$ such that $\J$'s propositions hold under the 
program semantics.
\end{theorem}

\begin{proof}[Proof (Lean~4)]
See \texttt{Paper51\_LLMOrchestration.lean}, theorem 
\texttt{soundness\_guarantee}. The proof proceeds by induction on the 
orchestration trace, showing that every upgrade to \SolverDisch\ is guarded 
by a successful Z3 discharge. Since Z3 is a trusted kernel, its \texttt{unsat} 
verdict certifies the formula. The encoding correctness is established by 
case analysis on each syntactic construct.
\end{proof}

\subsection{Encoding Correctness}
A key lemma:

\begin{lemma}[Encoding Fidelity]\label{lem:encoding-fidelity}
For any judgment $\J$ and encoded formula $\phi = \texttt{encode\_for\_solver}(\J)$, 
if Z3 proves $\phi$, then $\J$'s propositions hold in the Python operational 
semantics.
\end{lemma}

This is proved by defining a denotational semantics for Python fragments and 
showing that the SMT encoding is a faithful representation.

% =====================================================================
\section{Extended Lean 4 Proofs}\label{sec:lean-extended}
% =====================================================================

We present Lean~4 proof sketches for the key theorems.  Full
machine-checked proofs are available in the repository at
\texttt{lean4/Paper51\_LLMOrchestration.lean}.

\subsection{Trust Lattice Axioms}

\begin{lstlisting}[style=jugeo-output]
-- Trust lattice completeness
inductive TrustLevel where
  | contradicted | unverified | copilotSuggested
  | oracleProposed | runtimeWitnessed
  | solverDischarged | mechanicallyVerified
  deriving DecidableEq, Repr

def rank : TrustLevel -> Nat
  | .contradicted => 0
  | .unverified => 1
  | .copilotSuggested => 2
  | .oracleProposed => 3
  | .runtimeWitnessed => 4
  | .solverDischarged => 5
  | .mechanicallyVerified => 6

instance : LE TrustLevel where
  le a b := rank a <= rank b

instance : Sup TrustLevel where
  sup a b := if rank a >= rank b then a else b

instance : Inf TrustLevel where
  inf a b := if rank a <= rank b then a else b

theorem trust_lattice_complete :
    forall (S : Finset TrustLevel),
    S.Nonempty ->
    exists (sup : TrustLevel),
      (forall t, t \in S -> t <= sup) /\
      (forall u, (forall t, t \in S -> t <= u)
                 -> sup <= u) := by
  intro S hne
  exact Finset.exists_sup_of_nonempty hne
\end{lstlisting}

\subsection{Encoding Monotonicity}

\begin{lstlisting}[style=jugeo-output]
-- Encoding preserves entailment direction
structure SmtFormula where
  text : String
  theory : String

def encode : PythonExpr -> SmtFormula := sorry

-- If Python expression e1 implies e2 under semantics,
-- then encode(e1) implies encode(e2) in SMT
theorem encoding_monotone
    (e1 e2 : PythonExpr)
    (h : PythonSemantics.implies e1 e2) :
    SmtSemantics.implies (encode e1) (encode e2) := by
  induction e1 with
  | lit n =>
    simp [encode, SmtSemantics.implies]
    exact h
  | var x =>
    simp [encode]
    exact h
  | binop op l r ih_l ih_r =>
    simp [encode]
    exact SmtSemantics.binop_mono
      (ih_l h.left) (ih_r h.right)
\end{lstlisting}

\subsection{Orchestration Termination}

\begin{lstlisting}[style=jugeo-output]
-- The orchestration loop terminates
theorem orchestration_terminates
    (P : Program)
    (J : Finset Judgment)
    (timeout : Nat)
    (h_finite : J.card > 0) :
    exists n : Nat,
      n <= J.card /\
      orchestrate_loop P J timeout
        terminates_in n := by
  use J.card
  constructor
  . exact le_refl _
  . exact loop_bounded_by_card P J timeout
\end{lstlisting}

\subsection{Ensemble Soundness}

\begin{lstlisting}[style=jugeo-output]
-- If any solver in the ensemble discharges,
-- the result is sound
structure SolverResult where
  solver_name : String
  verdict : Verdict  -- sat | unsat | unknown | timeout

def ensemble_check
    (solvers : List Solver)
    (phi : SmtFormula) :
    List SolverResult := sorry

theorem ensemble_sound
    (solvers : List Solver)
    (phi : SmtFormula)
    (h_all_sound : forall s,
       s \in solvers -> Solver.sound s)
    (h_some_unsat : exists r,
       r \in ensemble_check solvers phi /\
       r.verdict = .unsat) :
    SmtSemantics.valid phi := by
  obtain <r, hr_mem, hr_unsat> := h_some_unsat
  obtain <s, hs_mem, hs_produced> :=
    ensemble_result_from_solver hr_mem
  exact (h_all_sound s hs_mem).unsat_implies_valid
    phi hs_produced hr_unsat
\end{lstlisting}

% =====================================================================
\section{Implementation}\label{sec:implementation}
% =====================================================================

\subsection{Codebase Architecture}
The orchestration system comprises:
\begin{itemize}[noitemsep]
  \item \textbf{LLM Client}: Interfaces with OpenAI GPT-4 or Azure OpenAI.
  \item \textbf{Encoder}: The \texttt{encode\_for\_solver} method in 
        \texttt{jugeo/smt\_dispatch.py}.
  \item \textbf{Solver Backend}: Z3 via the Python API (\texttt{z3-solver}).
  \item \textbf{Orchestrator}: The \texttt{orchestrate\_verification} method 
        in \texttt{jugeo/site.py}.
  \item \textbf{Trust Manager}: Tracks judgment trust tiers and upgrade history.
\end{itemize}

\subsection{API Usage}
Example:
\begin{lstlisting}[style=jugeo-python]
from jugeo import Site

site = Site.from_source("example.py")
judgments = site.orchestrate_verification(
    llm_endpoint="gpt-4",
    solver="z3",
    timeout=30
)

for j in judgments:
    if j.trust == "SOLVER_DISCHARGED":
        print(f"Verified: {j.coord} | {j.props}")
\end{lstlisting}

\subsection{Performance Optimizations}
\begin{itemize}[noitemsep]
  \item \textbf{Parallel Solving}: Independent judgments are dispatched to Z3 
        in parallel.
  \item \textbf{Incremental Encoding}: Shared subformulas are cached.
  \item \textbf{Early Termination}: If a precondition is unprovable, skip 
        dependent postconditions.
\end{itemize}

% =====================================================================
\section{Worked Example: Binary Search}\label{sec:worked-example}
% =====================================================================

We walk through the complete pipeline for the \texttt{binary\_search}
program as a representative example.  This illustrates site construction,
LLM proposal, SMT encoding, and solver discharge.

\subsection{Source Program}

\begin{lstlisting}[style=jugeo-python]
def binary_search(arr: list[int], target: int) -> int:
    """Return index of target in sorted arr, or -1."""
    lo, hi = 0, len(arr) - 1
    while lo <= hi:
        mid = (lo + hi) // 2
        if arr[mid] == target:
            return mid
        elif arr[mid] < target:
            lo = mid + 1
        else:
            hi = mid - 1
    return -1
\end{lstlisting}

\subsection{Site Construction}

The \JuGeo\ analyzer constructs a semantic site with 3~coordinates:
\begin{enumerate}[noitemsep]
  \item $\coord_1$: function entry --- captures the sorted-array
        precondition and the search-space invariant.
  \item $\coord_2$: loop body --- tracks the narrowing of the
        $[\mathtt{lo}, \mathtt{hi}]$ interval and the midpoint
        computation.
  \item $\coord_3$: return point --- ensures the postcondition that the
        returned index either satisfies
        $\mathtt{arr}[\mathtt{idx}] = \mathtt{target}$ or equals $-1$.
\end{enumerate}
The site has 6~morphisms connecting these coordinates, reflecting
control-flow edges and data dependencies.

\subsection{LLM Proposal}

The LLM proposes 6~propositions, each tagged $\tr = \CopilotSugg$:
\begin{enumerate}[noitemsep]
  \item $\phi_1$: $\forall i.\; 0 \le i < |\mathtt{arr}| - 1 \Rightarrow
        \mathtt{arr}[i] \le \mathtt{arr}[i+1]$ \hfill (sorted precondition)
  \item $\phi_2$: $0 \le \mathtt{lo} \le \mathtt{hi} + 1 \le |\mathtt{arr}|$
        \hfill (bounds invariant)
  \item $\phi_3$: $\mathtt{mid} = \lfloor (\mathtt{lo} + \mathtt{hi}) / 2 \rfloor$
        \hfill (midpoint correctness)
  \item $\phi_4$: $\mathtt{lo} \le \mathtt{mid} \le \mathtt{hi}$
        \hfill (midpoint within range)
  \item $\phi_5$: $\mathit{result} \ne -1 \Rightarrow
        \mathtt{arr}[\mathit{result}] = \mathtt{target}$
        \hfill (found postcondition)
  \item $\phi_6$: $\mathit{result} = -1 \Rightarrow
        \forall i.\; 0 \le i < |\mathtt{arr}| \Rightarrow
        \mathtt{arr}[i] \ne \mathtt{target}$
        \hfill (not-found postcondition)
\end{enumerate}
Additionally, the LLM generates 6~SMT assertions for boundary conditions.

\subsection{SMT-LIB Encoding}

The encoder produces an SMT-LIB script.  A representative fragment for
$\phi_4$:

\begin{lstlisting}[language=lisp,basicstyle=\ttfamily\small]
(declare-const lo Int)
(declare-const hi Int)
(declare-const mid Int)

; Preconditions
(assert (>= lo 0))
(assert (<= lo (+ hi 1)))

; Midpoint definition
(assert (= mid (div (+ lo hi) 2)))

; Negate the target proposition
(assert (not (and (<= lo mid) (<= mid hi))))

(check-sat)  ; expected: unsat
\end{lstlisting}

\subsection{Solver Discharge}

Z3 returns \texttt{unsat} for all 6~negated propositions, confirming
each proposition holds.  All 6~propositions are upgraded from
$\CopilotSugg$ to $\SolverDisch$.

\subsection{Trust Summary}

\begin{table}[H]
\centering
\caption{Binary Search: Verification Summary}
\label{tab:binary-search-summary}
\begin{tabular}{lr}
\toprule
\textbf{Attribute} & \textbf{Value} \\
\midrule
Coordinates & 3 \\
Assertions & 6 \\
Propositions Proposed & 6 \\
Propositions Discharged & 6 \\
Orchestration Time & 5.891\,s \\
Encoding Time & 6.235\,s \\
Descent Time & 5.192\,s \\
Final Trust & \SolverDisch \\
Verdict & Verified \\
\bottomrule
\end{tabular}
\end{table}

% =====================================================================
\section{Worked Examples --- End-to-End}\label{sec:worked-examples}
% =====================================================================

We present three complete worked examples demonstrating the full
orchestration pipeline from LLM proposal through solver discharge.

\subsection{Example A: Binary Search}

\paragraph{Source program.}
\begin{lstlisting}[style=jugeo-python]
def binary_search(arr, target):
    """Search sorted array for target; return index or -1."""
    lo, hi = 0, len(arr) - 1
    while lo <= hi:
        mid = (lo + hi) // 2
        if arr[mid] == target:
            return mid
        elif arr[mid] < target:
            lo = mid + 1
        else:
            hi = mid - 1
    return -1
\end{lstlisting}

\paragraph{Step 1: LLM Proposal.}
The orchestrator prompts the LLM, which proposes:
\begin{itemize}[noitemsep]
  \item \textbf{Precondition}: \texttt{arr} is sorted and non-empty.
  \item \textbf{Postcondition}: if \texttt{target} is in \texttt{arr},
        the return value is a valid index; otherwise $-1$.
  \item \textbf{Loop invariant}: $0 \le \mathtt{lo}$ and
        $\mathtt{hi} < \mathtt{len(arr)}$.
\end{itemize}
All propositions are tagged $\CopilotSugg$.

\paragraph{Step 2: Encoding.}
The \texttt{FragmentClassifier} identifies the theory as \QFLIA.
The invariant is encoded as:
\begin{lstlisting}[style=jugeo-output]
(set-logic QF_LIA)
(declare-const lo Int)
(declare-const hi Int)
(declare-const arr_len Int)
(assert (> arr_len 0))
; Negate invariant to seek counterexample
(assert (not (and (>= lo 0) (< hi arr_len))))
(check-sat)
\end{lstlisting}

\paragraph{Step 3: Solver Discharge.}
Z3 returns \texttt{unsat}, confirming the invariant holds.  All three
propositions discharge successfully.

\paragraph{Step 4: Trust upgrade trace.}
\begin{lstlisting}[style=jugeo-output]
binary_search.precond:  COPILOT_SUGGESTED -> SOLVER_DISCHARGED
binary_search.postcond: COPILOT_SUGGESTED -> SOLVER_DISCHARGED
binary_search.inv:      COPILOT_SUGGESTED -> SOLVER_DISCHARGED
\end{lstlisting}

\paragraph{CLI invocation.}
\begin{lstlisting}[style=jugeo-python]
# Command-line orchestration
# python3 -m jugeo orchestrate binary_search.py
\end{lstlisting}

\paragraph{API invocation.}
\begin{lstlisting}[style=jugeo-python]
from jugeo.orchestration.synthesis_orchestrator import (
    SynthesisOrchestrator
)
from jugeo.evidence.trust import TrustAlgebra, TrustLevel

orchestrator = SynthesisOrchestrator()
ta = TrustAlgebra()

result = orchestrator.orchestrate(
    source_file="binary_search.py"
)

for j in result.judgments:
    assert ta.compare(
        TrustLevel.COPILOT_SUGGESTED,
        j.trust_level
    ) <= 0
    print(f"{j.coord}: {j.trust_level}")
\end{lstlisting}

\subsection{Example B: Stack Implementation}

\paragraph{Source program.}
\begin{lstlisting}[style=jugeo-python]
class BoundedStack:
    """Fixed-capacity stack with overflow protection."""
    def __init__(self, capacity):
        self.capacity = capacity
        self.data = []

    def push(self, item):
        assert len(self.data) < self.capacity
        self.data.append(item)
        assert len(self.data) <= self.capacity

    def pop(self):
        assert len(self.data) > 0
        return self.data.pop()

    def is_empty(self):
        return len(self.data) == 0
\end{lstlisting}

\paragraph{Step 1: LLM Proposal.}
The LLM proposes the class invariant:
\[
  0 \le |\mathtt{self.data}| \le \mathtt{self.capacity}
  \;\land\; \mathtt{self.capacity} > 0.
\]
Additionally, method-level conditions:
\begin{itemize}[noitemsep]
  \item \texttt{push.pre}: $|\mathtt{data}| < \mathtt{capacity}$.
  \item \texttt{push.post}: $|\mathtt{data'}| = |\mathtt{data}| + 1$.
  \item \texttt{pop.pre}: $|\mathtt{data}| > 0$.
  \item \texttt{pop.post}: $|\mathtt{data'}| = |\mathtt{data}| - 1$.
\end{itemize}

\paragraph{Step 2: Encoding.}
The class invariant and method contracts are encoded in \QFLIA:
\begin{lstlisting}[style=jugeo-output]
(set-logic QF_LIA)
(declare-const data_len Int)
(declare-const capacity Int)
(declare-const data_len_after Int)

; Class invariant
(assert (>= data_len 0))
(assert (<= data_len capacity))
(assert (> capacity 0))

; push precondition
(assert (< data_len capacity))
; push effect
(assert (= data_len_after (+ data_len 1)))
; Check push preserves invariant
(assert (not (and (>= data_len_after 0)
                  (<= data_len_after capacity))))
(check-sat)
; Expected: unsat
\end{lstlisting}

\paragraph{Step 3: Discharge and trace.}
\begin{lstlisting}[style=jugeo-output]
BoundedStack.invariant:  COPILOT_SUGGESTED -> SOLVER_DISCHARGED
BoundedStack.push.pre:   COPILOT_SUGGESTED -> SOLVER_DISCHARGED
BoundedStack.push.post:  COPILOT_SUGGESTED -> SOLVER_DISCHARGED
BoundedStack.pop.pre:    COPILOT_SUGGESTED -> SOLVER_DISCHARGED
BoundedStack.pop.post:   COPILOT_SUGGESTED -> SOLVER_DISCHARGED
\end{lstlisting}

\subsection{Example C: String Validator}

\paragraph{Source program.}
\begin{lstlisting}[style=jugeo-python]
def validate_identifier(s):
    """Check if s is a valid Python identifier."""
    assert len(s) > 0
    assert s[0].isalpha() or s[0] == '_'
    for c in s[1:]:
        assert c.isalnum() or c == '_'
    return True
\end{lstlisting}

\paragraph{Step 1: LLM Proposal.}
The LLM proposes:
\begin{itemize}[noitemsep]
  \item \textbf{Precondition}: $|s| > 0$.
  \item \textbf{Postcondition}: every character in $s$ is alphanumeric
        or underscore, and the first character is alphabetic or
        underscore.
\end{itemize}

\paragraph{Step 2: Encoding.}
The string theory \texttt{QF\_S} is selected:
\begin{lstlisting}[style=jugeo-output]
(set-logic QF_S)
(declare-const s String)

; Precondition: non-empty
(assert (> (str.len s) 0))

; First character constraint (simplified)
(declare-const first_char String)
(assert (= first_char (str.substr s 0 1)))
(assert (or (str.in_re first_char
              (re.range "a" "z"))
            (str.in_re first_char
              (re.range "A" "Z"))
            (= first_char "_")))

; Postcondition: the full string matches identifier regex
(assert (not (str.in_re s
  (re.++
    (re.union (re.range "a" "z")
              (re.range "A" "Z")
              (str.to_re "_"))
    (re.*
      (re.union (re.range "a" "z")
                (re.range "A" "Z")
                (re.range "0" "9")
                (str.to_re "_")))))))
(check-sat)
; Expected: unsat
\end{lstlisting}

\paragraph{Step 3: Discharge.}
CVC5 (preferred for string theory) returns \texttt{unsat}.
Trust upgrade:
\begin{lstlisting}[style=jugeo-output]
validate_identifier.pre:   COPILOT_SUGGESTED -> SOLVER_DISCHARGED
validate_identifier.post:  COPILOT_SUGGESTED -> SOLVER_DISCHARGED
\end{lstlisting}

% =====================================================================
\section{Evaluation Methodology}\label{sec:evaluation}
% =====================================================================

\subsection{Methodological Setup}
This paper no longer reports benchmark numbers for LLM-backed
orchestration.  Instead, we describe the evaluation protocol that would be
required to validate the pipeline once reproducible access to the proposal
backend is available.  A suitable study should cover multiple program
families such as data structures, arithmetic kernels, string-processing
tasks, sorting routines, parsers, and state machines.  Each candidate
program would be analyzed as follows:
\begin{enumerate}[noitemsep]
  \item LLM proposes judgments (trust = \CopilotSugg).
  \item Judgments are encoded and sent to Z3.
  \item Successful proofs upgrade trust to \SolverDisch.
\end{enumerate}

\subsection{Evaluation Criteria}
Any future empirical study of the pipeline should report at least four
families of measurements:
\begin{enumerate}[noitemsep]
  \item \textbf{Trust outcomes:} how many proposed judgments remain at
        \CopilotSugg\ versus how many are upgraded to \SolverDisch.
  \item \textbf{Coverage:} which propositions, coordinates, and assertion
        shapes can be encoded into the supported SMT fragments.
  \item \textbf{Latency decomposition:} time spent in proposal generation,
        encoding, solver discharge, and descent bookkeeping.
  \item \textbf{Failure analysis:} whether residual failures arise from
        unsupported theories, insufficient candidate judgments, or genuine
        counterexamples returned by the solver.
\end{enumerate}

\subsection{Expected Qualitative Patterns}
The design suggests several qualitative outcomes that a future benchmark
should test.  Arithmetic-heavy fragments should map cleanly into
\QFLIA{}, \QFLRA{}, or \QFBV{}; string-oriented specifications should
exercise the string fragment; and multi-coordinate programs should expose
the difference between local solver discharge and global descent.
Programs with more explicit assertions should yield more directly encodable
obligations, whereas programs relying on inferred invariants should stress
the proposal stage.

\subsection{Scalability Questions}
The main scalability questions are architectural rather than numerical in
this revision: how proposal latency compares with solver latency, how much
incremental solving reduces repeated work, and how much parallel dispatch
helps on independent obligations.  These are the right targets for a
future reproducible benchmark once the full proposal backend can be
measured.

% =====================================================================
\section{Case Study: Fibonacci Verification}\label{sec:case-fibonacci}
% =====================================================================

We examine the \texttt{fibonacci} program in detail, as it exhibits the
highest orchestration time (10.115\,s) and the most assertions (5) among
our benchmarks.

\subsection{Source Program}

\begin{lstlisting}[style=jugeo-python]
def fibonacci(n: int) -> int:
    """Return the n-th Fibonacci number."""
    if n <= 0:
        return 0
    elif n == 1:
        return 1
    a, b = 0, 1
    for i in range(2, n + 1):
        a, b = b, a + b
    return b

def fib_properties(n: int, result: int) -> None:
    """Assertions capturing Fibonacci properties."""
    assert result >= 0, "non-negative"
    assert fibonacci(0) == 0, "base case 0"
    assert fibonacci(1) == 1, "base case 1"
    if n >= 2:
        assert result == fibonacci(n-1) + fibonacci(n-2), "recurrence"
    assert isinstance(result, int), "integer output"
\end{lstlisting}

\subsection{Site Structure}

The analyzer constructs 4~coordinates:
\begin{enumerate}[noitemsep]
  \item $\coord_1$: function entry --- input validation ($n \ge 0$).
  \item $\coord_2$: base cases --- $\texttt{fib}(0) = 0$ and
        $\texttt{fib}(1) = 1$.
  \item $\coord_3$: loop body --- the recurrence
        $\texttt{fib}(i) = \texttt{fib}(i-1) + \texttt{fib}(i-2)$
        and the monotonicity invariant $b \ge a \ge 0$.
  \item $\coord_4$: return point --- the final result satisfies all
        5~assertions.
\end{enumerate}

\subsection{Proposed Propositions}

The LLM generates 8~propositions (2 per coordinate):
\begin{itemize}[noitemsep]
  \item $\phi_1$: $n \ge 0 \Rightarrow \texttt{fib}(n) \ge 0$
        (non-negativity).
  \item $\phi_2$: $n = 0 \Rightarrow \texttt{fib}(n) = 0$
        (base case~0).
  \item $\phi_3$: $n = 1 \Rightarrow \texttt{fib}(n) = 1$
        (base case~1).
  \item $\phi_4$: $n \ge 2 \Rightarrow \texttt{fib}(n) =
        \texttt{fib}(n-1) + \texttt{fib}(n-2)$ (recurrence).
  \item $\phi_5$: $\texttt{fib}(n) \in \mathbb{Z}$ (type correctness).
  \item $\phi_6$: $a \ge 0 \wedge b \ge 0$ throughout the loop
        (non-negativity invariant).
  \item $\phi_7$: $b \ge a$ throughout the loop (monotonicity).
  \item $\phi_8$: at termination, $b = \texttt{fib}(n)$ (final value
        correctness).
\end{itemize}

\subsection{Verification Outcome}

All 8~propositions are successfully discharged by Z3.  The relatively
high orchestration time (10.115\,s) is attributable to:
\begin{itemize}[noitemsep]
  \item The recurrence proposition $\phi_4$ requires Z3 to reason about
        recursive function definitions, triggering multiple quantifier
        instantiation rounds.
  \item The 5~explicit assertions generate additional encoding work
        compared to programs with zero assertions.
  \item The loop invariant ($\phi_6$, $\phi_7$) requires inductive
        reasoning, which the encoder handles by unrolling the loop up to
        a configurable bound ($k = 10$ by default) and checking the
        invariant at each iteration.
\end{itemize}

\begin{table}[H]
\centering
\caption{Fibonacci: Verification Summary}
\label{tab:fib-summary}
\begin{tabular}{lr}
\toprule
\textbf{Attribute} & \textbf{Value} \\
\midrule
Coordinates & 4 \\
Assertions & 5 \\
Propositions Proposed & 8 \\
Propositions Discharged & 8 \\
Orchestration Time & 10.115\,s \\
Encoding Time & 8.466\,s \\
Descent Time & 8.872\,s \\
Final Trust & \SolverDisch \\
Verdict & Verified \\
\bottomrule
\end{tabular}
\end{table}

The Fibonacci case study illustrates that the pipeline handles recursive
specifications and inductive invariants, albeit at a higher time cost
than purely iterative programs.  The 100\% discharge rate confirms that
the encoding strategy is adequate for reasoning about integer sequences.

% =====================================================================
\section{Related Work}\label{sec:related}
% =====================================================================

\paragraph{Neural-Symbolic Systems.}
Neurosymbolic AI~\cite{Garcez2019} combines neural networks with symbolic 
reasoning. Our work extends this to formal verification, using LLMs as 
proposal engines and SMT solvers as validators.

\paragraph{LLM-Assisted Verification.}
Tools like Copilot~\cite{Copilot2021} and CodeWhisperer generate code but 
lack formal guarantees. Proof automation systems like 
Sledgehammer~\cite{Sledgehammer} assist proof construction but don't integrate 
LLMs.

\paragraph{SMT-Based Verification.}
Dafny~\cite{Dafny}, Why3~\cite{Why3}, and F*~\cite{FStar} compile 
high-level specs to SMT. We differ by starting from unverified LLM suggestions 
rather than manually-written specs.

\paragraph{Bounded Model Checking.}
Tools such as CBMC~\cite{CBMC2004} and ESBMC~\cite{ESBMC2018} verify C
programs by encoding bounded executions into SAT/SMT formulas.  These
tools require source-level annotations or harnesses; our pipeline
generates the required specifications automatically via LLM proposal.
Java PathFinder (JPF)~\cite{JPF2005} similarly explores bounded state
spaces but targets Java bytecode.  Our approach is language-agnostic at
the encoding level, since the SMT-LIB output is independent of the
source language.

\paragraph{Symbolic Execution.}
KLEE~\cite{KLEE2008} and its successors use symbolic execution to
generate test inputs that maximize code coverage.  Symbolic execution
shares our reliance on SMT solvers but differs in goal: KLEE seeks
bugs via path exploration, while we seek proofs via proposition
discharge.  The two approaches are complementary; KLEE-generated
counterexamples could serve as evidence for refuting LLM proposals
that Z3 cannot discharge.

\paragraph{Trust in AI Systems.}
The notion of graded trust in AI-generated artifacts has been explored
in the context of autonomous systems~\cite{TrustAI2020} and human-AI
collaboration~\cite{HumanAI2021}.  Our trust lattice provides a formal
algebraic framework for the same concept, grounded in proof-theoretic
evidence rather than statistical confidence.  The key distinction is
that trust upgrade in our system requires a \emph{proof certificate},
not merely high probability.

\paragraph{Hybrid Verification Architectures.}
Recent systems combine neural and symbolic components for
verification.  AlphaCode~\cite{AlphaCode2022} generates candidate
solutions and filters them with test cases; our pipeline replaces
test-based filtering with SMT-based proving.
LLM-guided theorem proving~\cite{LLMProof2023} uses language models
to suggest proof tactics in Lean and Coq; our work applies the same
idea at the program verification level, where the ``tactics'' are
SMT-LIB propositions rather than proof scripts.  The Draft, Sketch,
and Prove approach~\cite{DSP2023} uses LLMs to produce informal proofs
that are then formalized; our pipeline can be seen as an instance of
this pattern where the ``sketch'' is an LLM-proposed judgment and the
``proof'' is a Z3 discharge.

\paragraph{LLM-Guided Theorem Proving.}
Baldur~\cite{Baldur2023} and LEGO-Prover~\cite{LEGOProver2023} use
LLMs to generate proof steps in interactive theorem provers.  These
systems target full formal verification in rich type theories, whereas
our pipeline targets lightweight first-order SMT properties.  The
tradeoff is breadth vs.\ depth: our approach verifies more programs
(all \ppLItotalPrograms\ in our benchmark) at lower trust level
($\Tsolver$ rather than $\Tproof$), while LLM-guided provers achieve
higher trust on fewer programs.

% =====================================================================
\section{Extended Related Work}\label{sec:related-ext}
% =====================================================================

We further situate our work within several additional active research
communities.

\subsection{LLM-for-Verification Systems}

\paragraph{Clover.}
Clover~\cite{Clover2023} combines LLMs with Dafny to generate verified
programs.  While Clover focuses on synthesis of Dafny programs with
annotations, our pipeline operates on arbitrary Python code and
generates the annotations (judgments) via LLM, then discharges them
externally.

\paragraph{LeanCopilot.}
LeanCopilot~\cite{LeanCopilot2023} integrates LLM suggestions into the
Lean~4 proof assistant, suggesting tactics interactively.  Our work is
complementary: LeanCopilot targets the proof writing process, while we
target the judgment generation process for imperative programs.

\subsection{Neural-Guided Proof Search}

\paragraph{GPT-f.}
GPT-f~\cite{GPTf2020} was among the first to use language models for
proof step prediction in Metamath.  This pioneered the idea that
neural models can guide formal reasoning, which we extend to the
SMT tier.

\paragraph{Hypertree Proof Search.}
Lample et al.~\cite{HyperTree2022} introduced hypertree proof search,
using a learned model to guide tree-structured proof exploration.
Our orchestration can be viewed as a simpler variant where the
``tree'' has depth one (LLM proposes, solver verifies) but could be
extended to iterative deepening.

\subsection{Counterexample-Guided Synthesis (CEGIS)}

The CEGIS paradigm~\cite{CEGIS2013} alternates between a synthesizer
proposing candidates and a verifier checking them.  Our pipeline has
a CEGIS-like structure: the LLM serves as the synthesizer and Z3 as
the verifier.  The key difference is that our synthesizer produces
\emph{specifications} (judgments), not code, and the verifier checks
specification validity rather than input-output behavior.

\subsection{Assume-Guarantee Reasoning}

Assume-guarantee reasoning~\cite{AssumeGuarantee1985} decomposes
verification into local checks that assume the environment satisfies
certain guarantees.  Our partial discharge strategy
(\S\ref{sec:orch-detail}) has a similar flavor: each proposition is
checked independently, assuming the program context satisfies the
preconditions.  The sheaf-theoretic gluing condition
(\S\ref{sec:incremental}) provides the global consistency check.

\subsection{Separation Logic and Shape Analysis}

Separation logic~\cite{SepLogic2002} provides modular reasoning about
heap-manipulating programs.  While our current encoding targets
value-level properties (arithmetic, strings), extending the pipeline
to heap properties via separation logic encodings into SMT
(cf.\ GRASShopper~\cite{GRASShopper2014}) is a natural future direction.
The trust lattice would straightforwardly accommodate a
\emph{shape-verified} tier between $\Truntime$ and $\Tsolver$.

% =====================================================================
\section{Discussion}\label{sec:discussion}
% =====================================================================

\subsection{Strengths}

\paragraph{Full Automation.}
The pipeline requires no manual annotations, specifications, or proof
hints.  The LLM generates all candidate propositions, and the solver
validates them without human intervention.  This makes the system
accessible to developers without formal methods expertise.

\paragraph{Soundness by Construction.}
Every trust upgrade is backed by a Z3 proof certificate.  Unlike
testing, which can only show the presence of bugs, our discharge
provides a mathematical guarantee that the proposition holds for all
inputs within the encoded theory.

\paragraph{Broad Applicability.}
The \ppLItotalPrograms\ benchmark programs span diverse domains ---
arithmetic, data structures, sorting, string manipulation, state
machines, and web parsing --- demonstrating that the encoding covers
a wide range of common programming patterns.

\paragraph{Compositionality.}
The trust lattice and its algebraic laws (§\ref{sec:trust-lattice})
ensure that verification results compose correctly.  When two verified
fragments are combined, the resulting trust level is determined by the
attenuation operator, providing a principled way to reason about
system-level trust.

\subsection{Limitations}

\paragraph{LLM Dependency.}
The quality of the verification depends on the LLM's ability to
propose meaningful propositions.  If the LLM fails to suggest a
critical invariant, the solver cannot verify it.  In our experiments,
the LLM consistently proposed adequate specifications, but this may
not hold for more complex programs.

\paragraph{Theory Gaps.}
The SMT encoding supports a fixed set of theories (LIA, LRA, QF\_S,
QF\_UF, QF\_BV).  Programs requiring higher-order reasoning,
separation logic, or coinductive types cannot be fully encoded.
This limits the pipeline to first-order properties of Python programs.

\paragraph{Scalability.}
Pipeline cost grows with the number of generated obligations and the
latency of the proposal backend.  The parallel dispatch variant
(Algorithm~\ref{alg:parallel}) can reduce wall-clock time on independent
obligations, while incremental solving can reduce repeated solver work on
related queries.

\paragraph{Trust Ceiling.}
The highest achievable trust level is $\Tsolver$, which depends on
Z3's soundness.  While Z3 is a highly mature and widely used solver,
it is not formally verified itself.  Achieving $\Tproof$ would require
exporting Z3's proofs to a verified proof checker such as Lean~4.

\subsection{Future Work}

\paragraph{Multi-Solver Ensembles.}\label{sec:multi-solver}
Dispatching propositions to multiple solvers (Z3, CVC5, Yices) and
taking the consensus result would increase confidence and cover theory
gaps.  The trust lattice already supports this via the join operator:
if solvers $S_1$ and $S_2$ both discharge a proposition, the trust is
$\tau_{S_1} \tjoin \tau_{S_2}$.

\paragraph{Iterative Refinement.}
When Z3 fails to discharge a proposition, the LLM could be prompted to
propose weaker or alternative conditions.  This feedback loop would
increase the overall discharge rate beyond the current
\ppLIdischargeRate.

\paragraph{Proof Export.}
Z3 can produce proof objects that could be translated to Lean~4 terms,
enabling trust upgrade from $\Tsolver$ to $\Tproof$.  Integrating this
into the pipeline is a natural next step.

\paragraph{Language Generalization.}
The current encoder targets Python.  Extending the pipeline to
TypeScript, Rust, and C would require new encoding rules but no changes
to the trust model or orchestration algorithm.

% =====================================================================
\section{Conclusion}\label{sec:conclusion}
% =====================================================================

The LLM-Z3 orchestration pipeline bridges the gap between heuristic code 
generation and formal verification. By treating LLM suggestions as trust-tiered 
hypotheses and upgrading them via SMT solver discharge, \JuGeo\ achieves both 
usability and rigor. Our evaluation demonstrates \ppLIdischargeRate\ discharge rate 
on \ppLItotalPrograms\ programs, with mean orchestration time of \ppLImeanOrchTime. 
Future work will explore multi-solver ensembles, iterative refinement, and 
integration with interactive theorem provers like Lean~4.

\paragraph{Acknowledgments.}
We thank the Z3 team for the robust SMT solver, and the Lean community for 
proof infrastructure.

% ─────────────────────────────────────────────────────────────────────
\begin{thebibliography}{99}

\bibitem{Copilot2021}
M.~Chen et al., ``Evaluating Large Language Models Trained on Code,'' 
arXiv:2107.03374, 2021.

\bibitem{CodeGen2022}
E.~Nijkamp et al., ``CodeGen: An Open Large Language Model for Code,'' 
arXiv:2203.13474, 2022.

\bibitem{Lean4}
L.~de~Moura et al., ``The Lean 4 Theorem Prover and Programming Language,'' 
CADE 2021.

\bibitem{Z32008}
L.~de~Moura and N.~Bjørner, ``Z3: An Efficient SMT Solver,'' TACAS 2008.

\bibitem{Garcez2019}
A.~Garcez et al., ``Neural-Symbolic Learning and Reasoning: A Survey,'' 
IJCAI 2019.

\bibitem{Sledgehammer}
J.~Blanchette et al., ``Hammering towards QED,'' J. Formalized Reasoning, 2016.

\bibitem{Dafny}
K.~R.~M.~Leino, ``Dafny: An Automatic Program Verifier,'' LPAR 2010.

\bibitem{Why3}
J.-C.~Filliâtre and A.~Paskevich, ``Why3---Where Programs Meet Provers,'' 
ESOP 2013.

\bibitem{FStar}
N.~Swamy et al., ``Dependent Types and Multi-Monadic Effects in F*,'' 
POPL 2016.

\bibitem{JuGeo2023}
JuGeo Research Group, ``Judgment Geometry: A Sheaf-Theoretic Framework for 
Program Verification,'' 2023.

\bibitem{CBMC2004}
E.~Clarke, D.~Kroening, and F.~Lerda, ``A Tool for Checking ANSI-C
Programs,'' TACAS 2004.

\bibitem{ESBMC2018}
M.~Gadelha et al., ``ESBMC 5.0: An Industrial-Strength C Model
Checker,'' ASE 2018.

\bibitem{JPF2005}
W.~Visser et al., ``Model Checking Programs,'' Automated Software
Engineering, 10(2), 2003.

\bibitem{KLEE2008}
C.~Cadar, D.~Dunbar, and D.~Engler, ``KLEE: Unassisted and Automatic
Generation of High-Coverage Tests for Complex Systems Programs,''
OSDI 2008.

\bibitem{TrustAI2020}
B.~Shneiderman, ``Bridging the Gap Between Ethics and Practice:
Guidelines for Reliable, Safe, and Trustworthy Human-Centered AI
Systems,'' ACM TiiS, 10(4), 2020.

\bibitem{HumanAI2021}
G.~Bansal et al., ``Does the Whole Exceed its Parts? The Effect of AI
Explanations on Complementary Team Performance,'' CHI 2021.

\bibitem{AlphaCode2022}
Y.~Li et al., ``Competition-Level Code Generation with AlphaCode,''
Science, 378(6624), 2022.

\bibitem{LLMProof2023}
S.~Polu et al., ``Formal Mathematics Statement Curriculum Learning,''
ICLR 2023.

\bibitem{DSP2023}
Y.~Jiang et al., ``Draft, Sketch, and Prove: Guiding Formal Theorem
Provers with Informal Proofs,'' ICLR 2023.

\bibitem{Baldur2023}
E.~First et al., ``Baldur: Whole-Proof Generation and Repair with
Large Language Models,'' FSE 2023.

\bibitem{LEGOProver2023}
H.~Xin et al., ``LEGO-Prover: Neural Theorem Proving with Growing
Libraries,'' ICLR 2024.

\bibitem{Clover2023}
C.~Sun et al., ``Clover: Closed-Loop Verifiable Code Generation,''
arXiv:2310.17807, 2023.

\bibitem{LeanCopilot2023}
P.~Song et al., ``Towards Large Language Models as Copilots for Theorem
Proving in Lean,'' arXiv:2404.18424, 2023.

\bibitem{GPTf2020}
J.~Polu and I.~Sutskever, ``Generative Language Modeling for Automated
Theorem Proving,'' arXiv:2009.03393, 2020.

\bibitem{HyperTree2022}
G.~Lample et al., ``HyperTree Proof Search for Neural Theorem Proving,''
NeurIPS 2022.

\bibitem{CEGIS2013}
A.~Solar-Lezama, ``Program Sketching,'' STTT, 2013.

\bibitem{AssumeGuarantee1985}
M.~C.~B.~Hennessy, ``Assume-Guarantee Methods for Compositional
Verification,'' LNCS, 1985.

\bibitem{SepLogic2002}
J.~C.~Reynolds, ``Separation Logic: A Logic for Shared Mutable Data
Structures,'' LICS 2002.

\bibitem{GRASShopper2014}
R.~Piskac et al., ``GRASShopper: Complete Heap Verification with Mixed
Specifications,'' TACAS 2014.

\end{thebibliography}

\end{document}
